% ===== ETape1 type de document =====
\documentclass[12pt, a4paper]{report}
%===== ETape2 langue et accents =====
\usepackage[utf8]{inputenc}
\usepackage[T1]{fontenc}
\usepackage[french]{babel}
%===== ETape3 mise en page =====
\usepackage[margin=2.5cm]{geometry}
\usepackage{setspace}
\onehalfspacing
%===== ETape4 formule mathematiques et police times =====
\usepackage{amsmath}
\usepackage[version=4]{mhchem}
\usepackage{newtxtext, newtxmath}
%===== ETape5 tableaux =====
\usepackage{enumitem}
\usepackage{booktabs}
\usepackage{array}
\usepackage{multirow}
\usepackage{longtable}
%===== ETape6 figure et legende =====
\usepackage{graphicx}
\usepackage{xcolor}
\usepackage{tikz}
\usetikzlibrary{shapes.geometric}
\usepackage{pdflscape}
\usepackage{adjustbox}

\newcommand{\imgstd}[1]{%
  \adjustbox{width=11cm, height=8cm, keepaspectratio=false}{\includegraphics{#1}}%
}

\usepackage{float}
\usepackage[labelfont=bf]{caption}
\addto\captionsfrench{\renewcommand{\tablename}{Tableau}}
\addto\captionsfrench{\renewcommand{\listfigurename}{LISTE DES FIGURES}}
\addto\captionsfrench{\renewcommand{\listtablename}{LISTE DES TABLEAUX}}
\captionsetup{labelsep=colon}
%===== ETape7 titres =====
\usepackage{titlesec}
\titleformat{\chapter}[display]{\bfseries\centering\fontsize{14}{17}\selectfont}{\chaptertitlename\ \thechapter}{10pt}{\fontsize{14}{17}\selectfont}
\titleformat{\section}{\bfseries\fontsize{14}{17}\selectfont}{\thesection}{1em}{}
\titleformat{\subsection}{\bfseries\fontsize{12}{14.5}\selectfont}{\thesubsection}{1em}{}
\titleformat{\subsubsection}{\bfseries\itshape\fontsize{12}{14.5}\selectfont}{\thesubsubsection}{1em}{}
\titlespacing*{\chapter}{0pt}{0pt}{25pt}
\titleformat{\paragraph}{\bfseries\itshape\fontsize{12}{14.5}\selectfont}{\theparagraph}{1em}{}
\titlespacing*{\paragraph}{0pt}{12pt}{6pt}
%===== ETape8 bibliographies =====
\usepackage{csquotes}
\usepackage[backend=biber, style=ieee]{biblatex}
\addbibresource{references.bib}
%===== ETape9 liens cliquables =====
\usepackage{hyperref}
%===== Étape 10 le corps du document =====
%\includeonly{chapitres/chap1/chap1,chapitres/chap2/chap2}
\setcounter{secnumdepth}{4}
\setcounter{tocdepth}{4}
\begin{document}
\begin{titlepage}
\thispagestyle{empty}
\centering

{\bfseries\small RÉPUBLIQUE DU SÉNÉGAL\par}
\vspace{2mm}
\begin{tikzpicture}
  \fill[green!50!black] (0,0) rectangle (0.8,0.55);
  \fill[yellow!90!orange] (0.8,0) rectangle (1.6,0.55);
  \fill[red!85!black] (1.6,0) rectangle (2.4,0.55);
  \node[star, star points=5, star point ratio=2.25, fill=green!50!black, inner sep=0pt, minimum size=0.35cm] at (1.2,0.275) {};
\end{tikzpicture}\par
{\small\itshape Un Peuple -- Un But -- Une Foi\par}
\vspace{2mm}
{\bfseries\small Ministère de l'Enseignement supérieur, de la Recherche et de l'Innovation\par}
\vspace{2mm}
\rule{3cm}{0.4pt}\par
\vspace{3mm}
{\color{brown!70!black}\large UNIVERSITÉ DU SINE SALOUM EL-HÂDJ IBRAHIMA NIASS\par}
\vspace{4mm}
\includegraphics[height=2.4cm]{images/logo_ussein.png}\par

\vspace{8mm}
{\bfseries\large UFR Sciences Fondamentales et de l'Ingénieur\par}
\vspace{2mm}
{\bfseries\large Département Mathématiques, Informatique\par}

\vspace{10mm}
{\bfseries\large Mémoire de fin de cycle\par}
\vspace{2mm}
{\bfseries En vue de l'obtention du Diplôme de Licence\par}
\vspace{2mm}
{\bfseries Spécialité : Agro TIC\par}

\vspace{10mm}
\begin{tikzpicture}
  \node[draw=blue!60!black, line width=1.5pt, rounded corners=8pt, fill=blue!8, text width=14cm, align=center, inner sep=12pt]
  {\bfseries\large Thème :  Optimisation de l’irrigation et de la fertilisation en sols salins par l’agriculture de précision : étude comparative avec les pratiques traditionnelles en cultures maraîchères à Kaolack};
\end{tikzpicture}

\vfill
\noindent
\begin{minipage}[t]{0.48\textwidth}
  \raggedright
  {\bfseries\underline{Réalisé par :}}\\[4mm]
  {\bfseries Birahim DIAGNE}\\[2mm]
  {\bfseries Mariama SANE}
\end{minipage}
\hfill
\begin{minipage}[t]{0.48\textwidth}
  \raggedleft
  {\bfseries\underline{Encadré par :}}\\[4mm]
  {\bfseries Dr. Babacar DIOP}\\[2mm]
  {\bfseries M. Mamadou L. DIEDHIOU}
\end{minipage}

\vspace{10mm}
\begin{tikzpicture}
  % aile gauche
  \filldraw[fill=white, draw=blue!50!black, line width=0.8pt]
    (-2.4,0.85) -- (-7,0.95) -- (-6.1,0.48) -- (-7,0) -- (-2.4,0.1) -- cycle;
  % aile droite
  \filldraw[fill=white, draw=blue!50!black, line width=0.8pt]
    (2.4,0.85) -- (7,0.95) -- (6.1,0.48) -- (7,0) -- (2.4,0.1) -- cycle;
  % plis gris
  \filldraw[fill=gray!40, draw=blue!50!black, line width=0.6pt]
    (-2.5,1.0) -- (-1.8,1.0) -- (-2.5,1.2) -- cycle;
  \filldraw[fill=gray!40, draw=blue!50!black, line width=0.6pt]
    (2.5,1.0) -- (1.8,1.0) -- (2.5,1.2) -- cycle;
  % cadre central
  \filldraw[fill=white, draw=blue!50!black, line width=0.8pt]
    (-2.6,-0.05) rectangle (2.6,1.0);
  \node at (0,0.475) {Année universitaire 2024--2025};
\end{tikzpicture}
\vspace{5mm}
\end{titlepage}
\thispagestyle{empty}

\begin{center}
    {\fontsize{32}{38}\selectfont\bfseries Dédicaces\par}
\end{center}

\vspace{0.4cm}

{\fontsize{12}{15}\selectfont
\setstretch{1.05}
\setlength{\parindent}{0pt}
\setlength{\parskip}{0.25cm}

Nous dédions ce mémoire à toutes les personnes dont l'amour, les prières, les encouragements et le soutien ont contribué à la réalisation de ce travail.

À nos parents biologiques, \textit{Serigne Cheikh DIAGNE et Sokhna Rockhaya DIAGNE, El Hadji SANE et Bintanding DIATTA}, pour leur amour, leur patience, leurs prières et la confiance qu'ils nous ont toujours accordée. Qu'Allah leur accorde une longue vie.

À nos \textit{frères et sœurs}, pour leurs conseils, leurs encouragements et leur soutien financier.

À nos encadrants \textit{Dr B. DIOP et M. Mamadou L DIEDHIOU} et à nos maîtres de stage \textit{M. Aboubacar KEBE et M. Cheikh B DIALLO}, pour leur engagement, leur rigueur, leurs conseils et les compétences agronomiques qu'ils nous ont transmises.

À nos camarades de la \textit{cinquième promotion AgroTIC}, pour l'entraide et les expériences partagées tout au long de ce parcours.

\textbf{Mariama SANE}

Je dédie particulièrement ce travail à mes parents adoptifs, \textit{Ansoumana DJIBA et Banna DIA}, qui ont pris soin de moi depuis mon enfance et financé mes études jusqu'au baccalauréat, ainsi qu'à ma grande sœur \textit{Binta SANE}, qui a pris en charge mes études supérieures, mon logement et mon alimentation. Sans vous, ce parcours n'aurait pas été possible. 
Je le dédie également à mes colocataires de la chambre 4 de la RS4, en témoignage de ma reconnaissance pour leur bienveillance, leur sympathie et leur attention à mon égard.

\textbf{Birahim DIAGNE}

Je dédie ce travail à mon homonyme \textit{Birane DIAGNE} et à mon camarade \textit{Mouhamed Moustapha DIALLO}, qui m'a généreusement prêté son ordinateur pendant mes examens et durant la rédaction de ce mémoire, lorsque le mien est tombé en panne. À mon aîné et frère \textit{Gora MBAYE}, pour son soutien ; à mon guide, \textit{Serigne Moustapha Ndiaye}, pour ses conseils et son accompagnement indéfectible ; ainsi qu'à ma famille de l'USSEIN, \textit{Daara Munawirus Suduur}, pour leur esprit collectif.

\begin{center}
\itshape
\textbf{\textit{Ce mémoire vous est respectueusement dédié.}}
\end{center}

}
\thispagestyle{empty}

\begin{center}
    {\fontsize{28}{34}\selectfont\bfseries Remerciements\par}
\end{center}

\vspace{0.3cm}

{\fontsize{12}{15}\selectfont
\setstretch{1.1}
\setlength{\parindent}{0pt}
\setlength{\parskip}{0.2cm}

\textbf{\textit{Louange à Allah, Seigneur de l'univers.}}

Nous rendons grâce à \textbf{Allah, le Tout-Puissant}, par la volonté de qui toute chose aboutit, pour la santé, la force et la patience qu'Il nous a accordées afin de mener ce travail à son terme.

Nos remerciements vont d'abord à nos deux encadrants, \textit{Dr Babacar DIOP}, encadreur interne, et \textit{M. Mamadou Lamine DIEDHIOU}, encadreur externe. Tous deux se sont montrés disponibles tout au long de ce travail. Leurs conseils pertinents, leur rigueur, leur engagement, leur patience et leur expertise ont été déterminants pour mener ce mémoire à son terme.

Nous adressons également notre gratitude au \textit{Pr Abdoulaye SENE}, directeur de l'UFR, au \textit{Dr Modou LO}, chef du Département de Mathématiques et Informatique, au \textit{Pr Pascal François Mbissane FAYE}, responsable de la licence Agro TIC, au \textit{Dr Moussa GUEYE}, au \textit{Dr Madiop DIOUF}, au \textit{Dr Pape Djiby MERGANE}, à \textit{M. Bassirou DIENE} ainsi qu'à nos enseignants vacataires, pour la qualité de la formation reçue tout au long de notre parcours.

Nous renouvelons nos remerciements au \textit{Pr Pascal F. Mb FAYE}, pour avoir mis à notre disposition le capteur 8-en-1 durant tout le cycle expérimental.

Nos remerciements vont aux \textit{membres du jury} pour l'honneur qu'ils nous font en examinant ce mémoire et pour leurs remarques, qui contribueront à l'améliorer.

Nous remercions \textit{M. Idrissa TINE}, notre aîné, pour ses conseils en agroécologie.

Nous remercions nos aînés, \textit{Thiécoura NDIAYE, Mbaye Modou LO, Moustapha DIAKHATE, Bassine MBENGUE, Aïssata BALDE et Rachy S SOW}, pour leur soutien et le partage de leur expérience. 

Nos remerciements vont à nos camarades de la \textit{cinquième promotion AgroTIC}, pour leur entraide, leur solidarité et leur bienveillance.

Nous tenons également à souligner la qualité de notre collaboration et le soutien mutuel au sein de notre binôme : dans les moments de doute, notre soutien réciproque nous a rappelé que, dans la recherche, \textbf{toute erreur est une expérience}.

Enfin, nous remercions nos \textit{familles}, et en particulier nos \textit{parents}, pour leurs prières, leur patience et leurs encouragements.

Que toutes les personnes que nous n'avons pas pu citer et qui ont compté dans notre parcours trouvent ici l'expression de notre sincère gratitude.

}
\chapter*{RÉSUMÉ}
\addcontentsline{toc}{chapter}{RÉSUMÉ}
La salinité des sols, très présente au Sine-Saloum, associée à des pratiques d'irrigation et de fertilisation empiriques, compromet la production maraîchère. Ce mémoire, intitulé \textbf{« Optimisation de l’irrigation et de la fertilisation en sols salins par l’agriculture de précision : étude comparative avec les pratiques traditionnelles en cultures maraîchères à Kaolack »}, évalue l'intérêt d'une approche de précision dans ce contexte.

L'expérimentation, effectuée sur le champ d'application Agro TIC de l'USSEIN, compare deux zones de trois plots de $2\,\mathrm{m^2}$ (navet, gombo et laitue). La zone test suit une approche intégrée : étude du sol, amendement ( biochar, PCA, compost et EGB), puis irrigation et fertilisation pilotées à l'aide d'un capteur 8-en-1, de données météorologiques et de la méthode FAO-56 Penman-Monteith \cite{allen_crop_1998-1}. La zone témoin suit la pratique traditionnelle, avec un arrosage fixe. Une application web progressive, AgriSens, conçue avec la méthode MERISE et l'architecture FIWARE, automatise les calculs et la visualisation des données.

Les résultats montrent sur la zone test, un sol plus meuble et plus perméable, sans efflorescence saline, une meilleure croissance (gombo : 41,3~cm contre 21~cm) et une production effective (navet : 4,8~kg sur $2\,\mathrm{m^2}$), alors que dans la zone témoin la production est limitée voire inexistante. Les limites tiennent à l' estimation du NPK à partir de l'EC, à l'absence d'identification des ions, à la faible efficacité des traitements biologiques et à la détection des maladies par simple observation à l'œil nu.

\vspace{1em}
\noindent \textbf{Mots-clés :} agriculture de précision, sols salins, irrigation, fertilisation, évapotranspiration, amendement organique, capteur 8-en-1, cultures maraîchères, AgriSens.
\chapter*{ABSTRACT}
\addcontentsline{toc}{chapter}{ABSTRACT}

Soil salinity, widespread in the Sine-Saloum region, combined with empirical irrigation and fertilization practices, compromises vegetable production. This thesis, entitled ``Optimization of irrigation and fertilization through precision agriculture in saline soils: a comparative study with traditional practice on vegetable crops in Kaolack'', assesses the value of a precision approach in this context.

The experiment, carried out on the Agro TIC experimental field of USSEIN, compares two zones of three 2~m\textsuperscript{2} plots each (turnip, okra and lettuce). The test zone follows an integrated approach: soil study, amendment (biochar, PCA, compost and EGB), then irrigation and fertilization driven by an 8-in-1 sensor, weather data and the FAO-56 Penman-Monteith method. The control zone follows traditional practice, with fixed watering. A progressive web application, AgriSens, designed with the MERISE method and the FIWARE architecture, automates the calculations and the visualization of the data.

The results show that, in the test zone, the soil is looser and more permeable, with no salt efflorescence, along with better growth (okra: 41.3~cm versus 21~cm) and effective production (turnip: 4.8~kg over 2~m\textsuperscript{2}), whereas in the control zone production is limited or even nonexistent. The limitations lie in the estimation of NPK from EC, the absence of ion identification, the low efficacy of biological treatments, and the detection of diseases by simple visual observation.

\vspace{1em}
\noindent \textbf{Keywords:} precision agriculture, saline soils, irrigation, fertilization, evapotranspiration, organic amendment, 8-in-1 sensor, vegetable crops, AgriSens.
\chapter*{LISTE DES ABRÉVIATIONS}
\addcontentsline{toc}{chapter}{LISTE DES ABRÉVIATIONS}

\begin{longtable}{p{3.5cm}p{11cm}}
    \textbf{AP} & Agriculture de précision \\
    \textbf{FAO} & Organisation des Nations unies pour l'alimentation et l'agriculture \\
    \textbf{USDA} & Département de l'Agriculture des États-Unis \\
    \textbf{MO} & Matière organique \\
    \textbf{GPS} & \textit{Global Positioning System} \\
    \textbf{GSM} & Système mondial de communications mobiles \\
    \textbf{IoT} & Internet des objets \\
    \textbf{API} & \textit{Application Programming Interface} \\
    \textbf{CA} & Coque d'arachide \\
    \textbf{EGB} & Engrais granulé bio \\
    \textbf{PCA} & Poudre de coque d'arachide \\
    \textbf{CEC} & Capacité d'échange cationique \\
    \textbf{HT-PGPR} & \textit{Halotolerant Plant Growth-Promoting Rhizobacteria} \\
    \textbf{EC} & ElectroConductivité ou Conductivité électrique \\
    \textbf{N} & Azote \\
    \textbf{P} & Phosphore \\
    \textbf{K} & Potassium \\
    \textbf{GOM T} & Gombo témoin \\
    \textbf{GOM TS} & Gombo test \\
    \textbf{LAI T} & Laitue témoin \\
    \textbf{LAI TS} & Laitue test \\
    \textbf{NAV T} & Navet témoin \\
    \textbf{NAV TS} & Navet test \\
    \textbf{Plot T} & Plot témoin \\
    \textbf{Plot TS} & Plot test \\
    \textbf{RDU} & Réserve difficilement utilisable \\
    \textbf{RFU} & Réserve facilement utilisable \\
    \textbf{RU} & Réserve utile \\
    \textbf{SAD} & Système d'aide à la décision \\
    \textbf{MERISE} & Méthode d'étude et de réalisation informatique pour les systèmes d'entreprise \\
    \textbf{MCC} & Modèle conceptuel de communication \\
    \textbf{MCT} & Modèle conceptuel de traitement \\
    \textbf{MOT} & Modèle organisationnel de traitement \\
    \textbf{MCD} & Modèle conceptuel des données \\
    \textbf{MLD} & Modèle logique des données \\
    \textbf{PWA} & \textit{Progressive Web App} \\
\end{longtable}

\chapter*{LISTE DES SYMBOLES}
\addcontentsline{toc}{chapter}{LISTE DES SYMBOLES}

\begin{longtable}{p{3.5cm}p{11cm}}
    $\Delta$ & Pente de la courbe de pression de vapeur \\
    $\Delta S$ & Variation de la réserve en eau \\
    $\delta$ & Déclinaison solaire \\
    $\gamma$ & Constante psychrométrique \\
    $\omega_s$ & Angle horaire du coucher du soleil \\
    $\theta_{\mathrm{CC}}$ ou $H_{cc}$ & Humidité à la capacité au champ \\
    $\theta_{\mathrm{PFP}}$ ou $H_{pf}$ & Humidité au point de flétrissement permanent \\
    $\theta_{\text{actuel}}$ & Teneur en eau volumique actuelle du sol \\
    $D$ & Drainage profond \\
    $D_i$ & Dose d'irrigation \\
    $d_r$ & Distance relative Terre-Soleil \\
    $DPV$ & Déficit de pression de vapeur \\
    $e_a$ & Pression de vapeur réelle \\
    $e_s$ & Pression de vapeur saturante \\
    $ET_0$ & Évapotranspiration de référence \\
    $ET_c$ & Évapotranspiration culturale \\
    $ET_{c,\mathrm{aj}}$ & Évapotranspiration culturale ajustée \\
    $G$ & Flux de chaleur dans le sol \\
    $I$ & Irrigation \\
    $J$ & Numéro du jour de l'année \\
    $K_c$ & Coefficient cultural \\
    $K_{c,\mathrm{ini}}$ & Coefficient cultural initial \\
    $K_{c,\mathrm{mid}}$ & Coefficient cultural de mi-saison \\
    $K_{c,\mathrm{end}}$ & Coefficient cultural de fin de saison \\
    $K_{cb}$ & Coefficient cultural de base \\
    $K_e$ & Coefficient d'évaporation du sol \\
    $K_{Rs}$ & Coefficient d'ajustement \\
    $K_s$ & Coefficient de stress hydrique \\
    $p$ & Fraction de déplétion \\
    $P$ & Pression atmosphérique (précipitation dans le bilan hydrique) \\
    $R$ & Ruissellement \\
    $R_a$ & Rayonnement extraterrestre \\
    $R_c$ & Remontée capillaire \\
    $R_n$ & Rayonnement net \\
    $R_{nl}$ & Rayonnement net de grandes longueurs d'onde \\
    $R_{ns}$ & Rayonnement net de courtes longueurs d'onde \\
    $R_s$ & Rayonnement solaire \\
    $R_{so}$ & Rayonnement de ciel clair \\
    $RH_{\max}$ & Humidité relative maximale \\
    $RH_{\min}$ & Humidité relative minimale \\
    $S_a$ & Stock d'eau actuel \\
    $S_c$ & Seuil critique \\
    $T$ & Température moyenne journalière \\
    $T_{\max}$ & Température maximale \\
    $T_{\min}$ & Température minimale \\
    $u_2$ & Vitesse du vent à 2 m de hauteur \\
    $Z_r$ & Profondeur racinaire \\
    $z$ & Altitude du site \\
\end{longtable}
\tableofcontents
\clearpage
\addcontentsline{toc}{chapter}{LISTE DES FIGURES}
\listoffigures
\clearpage
\addcontentsline{toc}{chapter}{LISTE DES TABLEAUX}
\listoftables
\chapter*{INTRODUCTION GÉNÉRALE}
\addcontentsline{toc}{chapter}{INTRODUCTION GÉNÉRALE}
L’agriculture fait face à plusieurs défis, notamment la gestion durable des ressources en eau, l’amélioration des rendements et la préservation de l’environnement. Dans ce contexte, l’agriculture de précision constitue une approche permettant d’adapter les pratiques agricoles aux besoins des cultures grâce aux technologies modernes, telles que les capteurs, les stations météorologiques et les outils informatiques. Elle s’appuie sur des données mesurées pour améliorer la prise de décision et optimiser les apports en eau et en intrants.

Dans les systèmes agricoles traditionnels, l’irrigation est souvent pratiquée à des horaires réguliers, sans tenir suffisamment compte des besoins réels des cultures ni des conditions du sol et du climat. Ces pratiques peuvent entraîner un gaspillage d’eau, une utilisation inadaptée des intrants et une dégradation des sols. Ces contraintes sont particulièrement préoccupantes dans le Sine-Saloum, où la salinité des sols constitue un obstacle à la production maraîchère.

C’est dans ce contexte que s’inscrit cette étude, menée sur le site AgroTIC de l’USSEIN à Kaolack. Elle vise à évaluer une approche expérimentale combinant l’étude et l’amendement du sol, le suivi des paramètres agronomiques et le pilotage des apports hydriques. Cette démarche présente des enjeux environnementaux, économiques et scientifiques, notamment la réduction du gaspillage d’eau, l’optimisation des apports en intrants et l’évaluation comparative de l’agriculture de précision par rapport aux pratiques traditionnelles.
Quatre hypothèses de recherche ont été formulées.

\begin{itemize}[label=$\bullet$]
   	 \item H1 : les plots expérimentaux présentent un rendement supérieur par rapport aux plots témoins grâce à l’amendement et à une gestion précise et adaptée de l’eau et des intrants.
     
     \item H2 : la méthode de pilotage de l’irrigation et de la fertilisation établie permet de réduire la consommation d’eau et d’intrants de manière significative, abaissant ainsi les coûts de production.

     \item H3 : le suivi en temps réel des paramètres agronomiques permet une détection précoce des anomalies et une meilleure résilience face aux stress climatiques et/ou hydriques.

     \item H4 : malgré sa complexité apparente, l’application PWA est adoptable par les agriculteurs locaux après une courte durée de formation.
	\end{itemize}

L’objectif général de cette étude est de mettre en évidence l’efficacité de l’agriculture de précision dans l’optimisation de l’irrigation et des apports en intrants à travers une étude comparative. Plus précisément, il s’agit de :

\begin{itemize}[label=$\bullet$]
   	\item Mettre en place une parcelle expérimentale divisée en deux zones, une zone test pour l’agriculture de précision et une zone témoin pour l’agriculture traditionnelle, chacune composée de trois plots de  $2\,\mathrm{m^2}$, chaque plot ayant une culture différente de celles des autres ;
     \item Collecter des données de sol sur les plots tests en temps réel à l’aide du capteur 8-en-1 et de données météorologiques ;
     \item Estimer les besoins en eau des cultures en tenant compte de l’évapotranspiration et des pertes par infiltration ; 
     \item Déterminer les besoins en intrants en fonction des paramètres mesurés du sol (NPK) ; 
     \item Comparer les résultats obtenus entre les plots tests et les plots témoins (consommation d’eau, croissance et développement des cultures, rendement) ;
     \item Concevoir une application web progressive (PWA) permettant d’automatiser les calculs, de visualiser les données et de faciliter la prise de décision.
\end{itemize}

Ce mémoire est organisé en cinq chapitres. Le premier chapitre présente la revue de la littérature. Le deuxième décrit le protocole d’expérimentation. Le troisième présente la mise en œuvre et le suivi de l’expérimentation. Le quatrième présente la modélisation, la conception et l’implémentation de l’application. Enfin, le cinquième expose les résultats obtenus et leurs discussions.

\chapter{REVUE DE LA LITTÉRATURE}
\section{Introduction}
Face à la raréfaction des ressources en eau, aux effets du changement climatique et à la dégradation des sols, notamment la salinisation
qui compromet la productivité de nombreuses terres agricoles au Sénégal, l’amélioration des pratiques agricoles constitue un défi primordial.
Dans ce contexte, l’agriculture de précision vise à optimiser la gestion des ressources grâce aux capteurs, aux données météorologiques, aux 
technologies numériques et aux outils d’aide à la décision. Elle permet d’adapter les apports en eau et en intrants aux besoins réels des cultures
plutôt que de recourir à des méthodes fixes et empiriques, contribuant ainsi à améliorer les rendements et à préserver les ressources hydriques.
\section{Agriculture de précision :  concepts fondamentaux}

\subsection{Définition et principes de l'agriculture de précision}
L'agriculture de précision est une stratégie de gestion agricole fondée sur la prise en compte de la variabilité spatiale et temporelle des parcelles afin d'améliorer la durabilité agricole \cite{monteiro_precision_2021}. Contrairement aux méthodes conventionnelles, qui appliquent des interventions uniformes sur l'ensemble d'une exploitation, elle adapte les apports en eau, en fertilisants et en amendements aux besoins réels des cultures, en fonction des caractéristiques du sol, des conditions climatiques et du stade de développement des plantes.

Cette méthode repose sur la collecte, le traitement et l'analyse de données issues de différentes technologies, notamment les capteurs de sol, les stations météorologiques, les systèmes de géolocalisation (GPS), la télédétection, les drones et les analyses agronomiques \cite{shaheb_precision_2022}. Les informations recueillies permettent d'évaluer l'état des cultures et de leur environnement, facilitant ainsi une prise de décision plus précise.

L'agriculture de précision s'articule autour de trois principes fondamentaux : observer, analyser et agir \cite{noauthor_precision_2026}. Il s'agit d'abord de décrire la variabilité des parcelles, puis analyser les données afin d'identifier les besoins des cultures et enfin, d'ajuster les interventions (l'irrigation et la fertilisation), en fonction des conditions réelles du champ. 
L'objectif est d’améliorer l'efficacité de la gestion des ressources, de réduire les impacts environnementaux et de maintenir, voire augmenter les rendements.

\subsection{Origines et évolution du concept}
Le concept d’agriculture de précision s’est développé progressivement pour mieux prendre en compte la variabilité des sols, des cultures et des conditions de  l'environnement sur une même parcelle. 
Initialement fondée sur des observations et des mesures des propriétés du sol et des rendements, cette méthode a évolué avec l’introduction des outils de cartographie, du GPS et des systèmes d’information géographique, permettant une gestion différenciée des parcelles.

Les progrès des capteurs, de la télédétection et de l’informatique ont ensuite facilité la collecte et l’analyse des données météorologiques et agronomiques. 
Aujourd’hui, l’intégration de technologies connectées, notamment les \textbf{réseaux de capteurs et l’Internet des objets (IoT)} favorise un suivi plus régulier des parcelles et un ajustement des interventions en fonction des besoins réels des cultures. 
Cette évolution vise à améliorer l’utilisation des ressources, à \textbf{réduire les pertes et à préserver la productivité agricole}.

\subsection{Agriculture de précision versus agriculture traditionnelle}
L'agriculture traditionnelle repose généralement sur la gestion uniforme des parcelles, avec des apports en eau et en intrants 
déterminés selon l'expérience de l'agriculteur, un calendrier fixe ou l'observation visuelle des cultures. 
En revanche, l'AP s'appuie sur des capteurs et des données météorologiques afin d'adapter les interventions face aux variations du sol, du climat et des besoins réels des cultures.

Elle vise à améliorer l'utilisation de l'eau et des intrants, à limiter les pertes et à maintenir, voire augmenter les rendements. 
Toutefois, son adoption nécessite des compétences techniques, des équipements et un investissement initial important. Elle cherche donc à améliorer les pratiques traditionnelles en les rendants plus éfficaces.

\subsection{Enjeux et intérêts pour l'agroécologie}
L'agriculture de précision contribue à l'agroécologie en optimisant l'utilisation des intrants, des amendements organiques et de l'eau selon les caractéristiques du sol et les besoins réels des cultures. 
Elle permet ainsi de réduire le gaspillage des ressources et les impacts environnementaux, tout en améliorant l'efficacité des pratiques agricoles. 
Dans les zones confrontées à la raréfaction de l'eau et à la dégradation des sols comme le Sine Saloum, elle constitue une méthode pertinente pour concilier productivité agricole, préservation des ressources naturelles et durabilité des systèmes de production.

\section{Technologies de collecte de données en agriculture de précision}

\subsection{Les capteurs de sol et de climat}
En agriculture de précision, les capteurs constituent des outils essentiels car ils permettent de mesurer les paramètres du sol et de l'environnement. Les capteurs de sol mesurent notamment la température, l'humidité, la EC, le pH ainsi que les teneurs en N, P, K. 
Les appareils combinés, comme les capteur 8-en-1 regroupent plusieurs fonctions dans un même dispositif facilitant la collecte de données.

Les données météorologiques telles que l'humidité de l'air, la température et la vitesse du vent permettent d'estimer les besoins en eau des plantes. 
L'association des données du climat et du sol contribue ainsi à une meilleure gestion de l'irrigation et de la fertilisation.

\begin{figure}[H]
    \centering
    \imgstd{images/01_capteur_8_en_1.jpg}
    \caption{Capteur de sol 8-en-1. }
    \label{fig:8_en_1}
\end{figure}

\subsection{Les drones et l’imagerie aérienne en agriculture}
Les drones équipés de capteurs multispectraux ou optiques permettent de surveiller l'état des cultures et de détecter les zones de stress nutritif et/ou hydrique. Toutefois, leur utilisation nécessite des compétences techniques et des équipements spécifiques. Cependant cet outil n'a pas été utilisée dans notre expérimentation.

\subsection{Les donnees météorologiques et les API météo}
Les données climatiques telles que, la $T_{max}$ et $T_{min}$, la $RH_{max}$ et $RH_{min}$ ainsi que la vitesse du vent sont indispensables à l'estimation des besoins en eau des cultures. 
Pendant l'expérimentation, ces données ont été relevées à partir de l'application AccuWeather du smartphone. A la fin de l'expérimentation, l'API OpenWeather a été retenue pour automatiser la récupération des paramètres météorologiques de l'application conçue.

\subsection{Limites et défis de fiabilité des capteurs low-cost}
Malgré leur accessibilité, les \textbf{capteurs à faible coût} peuvent présenter des erreurs de mesure. La calibration, l'entretien et la vérification régulière des données sont donc nécessaires afin d'assurer leur fiabilité et d'améliorer la précision des décisions agricoles.


\section{Applications numériques et fondements théoriques de la gestion de l’irrigation}
\subsection{Panorama des applications numériques d’aide à la prise de décision d’irrigation}
Les applications numériques d'aide à la décision en irrigation permettent d'exploiter les données agronomiques, pédologiques et météorologiques pour optimiser la gestion de l'eau. 
Plusieurs outils et systèmes ont été développés avec de différentes approches et fonctionnalités.

{\small
\begin{longtable}{|p{3.4cm}|p{5.6cm}|p{5cm}|}
    \caption{Comparaison de quelques outils et travaux relatifs à la gestion numérique de l'irrigation}
    \label{tab:comparaison_outils_irrigation} \\
    \hline
    \textbf{Outil ou travail} & \textbf{Approche utilisée} & \textbf{Principaux apports} \\
    \hline
    \endfirsthead
    \hline
    \textbf{Outil ou travail} & \textbf{Approche utilisée} & \textbf{Principaux apports} \\
    \hline
    \endhead
    CROPWAT (FAO) & Calcul des besoins en eau des cultures et planification de l'irrigation à partir de données climatiques, culturales et pédologiques. & Aide à déterminer les besoins en eau et à planifier l'irrigation. \\
    \hline
    IrrigaSys & Système d'aide à la décision utilisant des données météorologiques et un modèle de bilan hydrique du sol. & Propose des calendriers d'irrigation accessibles sur plusieurs supports numériques. \\
    \hline
    Bluleaf & Application mobile intégrant des données météorologiques, culturales et des mesures d'humidité du sol. & Facilite la planification de l'irrigation à partir du bilan hydrique et des prévisions météorologiques. \\
    \hline
    Irrigator Pro & Système d'aide à la décision combinant les caractéristiques culturales, les données météorologiques et les mesures du sol. & Aide à programmer l'irrigation et a évolué vers une utilisation sur le web et les appareils mobiles. \\
    \hline
    idCROP & Application mobile associant des modèles de simulation des cultures, des données météorologiques et une analyse économique. & Propose des calendriers d'irrigation et permet d'évaluer les conséquences économiques des stratégies retenues. \\
    \hline
\end{longtable}
}

\begin{itemize}[label=$\bullet$]
    \item \textbf{Apport et positionnement de notre application AgriSens}
    
    Dans ce contexte, notre application se présente comme un instrument d’aide à la prise de décision hybride et accessible. Elle comble les insuffisances des solutions actuelles en associant : 
    \begin{itemize}[label=$\circ$]
    \item Modélisation hydrique dynamique : grâce au calcul en temps réel de $ET_0$ (avec la formule la plus recommandée) et de $ET_c$, en tenant compte du $K_c$ et du stade de développement de chacune de nos cultures choisies.
    
    \item Une évaluation complète de la santé du sol : en associant le contrôle de l’eau à la surveillance permanente des indicateurs physico-chimiques (pH, EC, N, P, K) mesurées du terrain.

    \item Une architecture et une ergonomie sur mesure : Une interface accessible sur mobile et web, offrant une centralisation multi-parcelles et un historique détaillé, garantissant une gestion rigoureuse et rationnelle de l’eau.
    \end{itemize}
\end{itemize}

\subsection{Approches de calcul des besoins en eau des cultures et de l'évapotranspiration}
Les besoins en eau d’une culture correspondent à la quantité d’eau nécessaire pour compenser les pertes principalement liées à l'évapotranspiration, qui regroupe \textbf{l'évaporation} de l'eau du sol et la \textbf{transpiration} des plantes. Trois grandeurs sont généralement distinguées : 

\begin{table}[H]
    \centering
    \caption{Correspondance entre grandeur, définition et utilisation}
    \label{tab:grandeurs_et}
    \begin{tabular}{|p{2.8cm}|p{6cm}|p{5.5cm}|}
        \hline
        \textbf{Grandeur} & \textbf{Définition} & \textbf{Utilisation} \\
        \hline
        $ET_0$ & Évapotranspiration d'une surface végétale de référence & Caractérisation de la demande climatique \\
        \hline
        $ET_c$ & Évapotranspiration d'une culture donnée & Calcul des besoins culturaux \\
        \hline
        $ET_{c,\mathrm{aj}}$ & Évapotranspiration en situation de stress ou de contraintes & Gestion en conditions réelles \\
        \hline
    \end{tabular}
\end{table}

\begin{itemize}[label=$\bullet$]
   	\item \textbf{Approche climatique}
 
    Cette approche estime l'évapotranspiration à partir des données météorologiques. La méthode FAO-56 de Penman-Monteith \cite{allen_crop_1998-1} constitue la référence recommandée pour calculer $ET_0$. D'autres méthodes, comme Hargreaves-Samani \cite{koc_reference_2023} peuvent être utilisées lorsque les données climatiques disponibles sont limitées.

   	\item \textbf{Approche par bilan hydrique du sol}
   	
    Le bilan hydrique permet de représenter les entrées et les sorties d’eau dans la zone racinaire sous une forme simplifiée :

    \begin{equation}
    \Delta S = P + I + R_c - ET_c - D - R
    \label{eq:bilan_hydrique}
    \end{equation}

    La $RU$ représente la quantité d’eau disponible entre la capacité au champ et le point de flétrissement permanent :
    \begin{equation}
    RU = 1000\,(\theta_{CC} - \theta_{PFP})\,Z_r
    \label{eq:ru}
    \end{equation}

    La $RFU$, correspondant à la fraction de la $RU$ que la culture peut consommer avant l’apparition d’un stress hydrique significatif, est donnée par :
    \begin{equation}
        RFU = RU \times p
        \label{eq:rfu}
    \end{equation}

    \item \textbf{Approche fondée sur les capteurs}
    
    Les capteurs d'humidité permettent le suivi de l'état hydrique du sol et le déclenchement de l'irrigation lorsque l'humidité descend sous un seuil défini. Le suivi à différentes profondeurs peut aussi aider à vérifier que l'eau atteint la zone racinaire sans provoquer de drainage excessif.

    \item \textbf{Approche directe et télédétection}

    Les méthodes directes, comme les lysimètres permettent de mesurer l'évapotranspiration, mais leur coût limite leur utilisation. La télédétection exploite les données aériennes ou satellitaires afin de suivre l'état hydrique et le développement des cultures sur des surfaces grandes. Ces méthodes n'ont pas été utilisées directement dans notre expérimentation, mais elles constituent des outils complémentaires.

\end{itemize}

\subsection{Méthode FAO-56 de Penman-Monteith, coefficient cultural (\texorpdfstring{$K_c$}{Kc}) et évapotranspiration culturale (\texorpdfstring{$ET_c$}{ETc})}

La méthode FAO-56 de Penman-Monteith \cite{allen_crop_1998-1} constitue la méthode de référence recommandée pour déterminer l'évapotranspiration de référence. Elle combine le bilan énergétique de la surface et les mécanismes aérodynamiques qui influencent le transfert de vapeur d’eau vers l'atmosphère. A l’echelle journaliere, l'équation s'écrit :
\begin{equation}
    ET_0 = \frac{0{,}408\,\Delta\,(R_n - G) + \gamma\,\dfrac{900}{T + 273}\,u_2\,(e_s - e_a)}{\Delta + \gamma\,(1 + 0{,}34\,u_2)}
    \label{eq:penman_monteith}
\end{equation}
Les principales variables de cette équation sont présentées dans le tableau ci-dessous.

\begin{table}[H]
    \centering
    \caption{Correspondance entre symboles et signification pour la formule $ET_0$}
    \label{tab:symboles_et0}
    \begin{tabular}{|p{3cm}|p{10cm}|}
        \hline
        \textbf{Symbole} & \textbf{Signification} \\
        \hline
        $R_n$ & Rayonnement net à la surface en MJ/m\textsuperscript{2}/jour \\
        \hline
        $G$ & Flux de chaleur dans le sol en MJ/m\textsuperscript{2}/jour \\
        \hline
        $T$ & Température moyenne de l'air à 2 m en °C \\
        \hline
        $u_2$ & Vitesse du vent à 2 m, en m/s \\
        \hline
        $e_s$ & Pression de la vapeur saturante, en kPa \\
        \hline
        $e_a$ & Pression de la vapeur réelle, en kPa \\
        \hline
        $e_s - e_a$ & Déficit de pression de vapeur, en kPa \\
        \hline
        $\Delta$ & Pente de la courbe de pression de vapeur, en kPa/°C \\
        \hline
        $\gamma$ & Constante psychrométrique, en kPa/°C \\
        \hline
    \end{tabular}
\end{table}

Le calcul de l’évapotranspiration culturale repose ensuite sur le coefficient cultural $K_c$, qui traduit les différences entre les besoins hydriques d’une culture et ceux de la surface de référence \cite{allen_crop_1998-1}. L’$ET_c$ est calculée par : 
\begin{equation}
    ET_c = K_c \times ET_0
    \label{eq:etc}
\end{equation}
Le $K_c$ varie selon le stade de développement de la culture. Le cycle cultural est généralement divisé en quatre phases : initiale, développement, mi-saison et fin de saison. Les valeurs caractéristiques sont notées $Kc_ini$, $Kc_mid$ et $Kc_end$.

Des ajustements peuvent être essentiels selon les conditions climatiques, la variété et les pratiques culturales. Lorsque la disponibilité en eau devient insuffisante, 
l'évapotranspiration réelle peut être inferieure à l'$ET_c$ standard. Un coefficient de stress hydrique ($K_s$) compris entre 0 et 1 peut alors être utilisé afin d'ajuster l'estimation des besoins en eau. 
Dans l'approche à coefficient cultural unique, cet ajustement s'écrit sous cette forme : 
\begin{equation}
    ET_{c,\mathrm{aj}} = (K_s \times K_{cb} \times ET_0) + (K_e \times ET_0)
    \label{eq:etc_aj}
\end{equation}
Les développements mathématiques complémentaires et les étapes détaillées de calcul des paramètres sont présentés en annexe (voir l’Annexe~\ref{ann:calculs}).

\subsection{Travaux antérieurs orientés sur différentes zones}

\begin{itemize}[label=$\bullet$]
   	\item \textbf{Travaux et orientations en zones tropicales}

En zones tropicales, l’irrigation numérique concerne notamment les cultures maraîchères, le riz, le maïs et les cultures sous serre. Les systèmes utilisent des capteurs, des microcontrôleurs, une transmission GSM ou LoRa et une interface mobile. Ces conditions soulèvent plusieurs difficultés particulières : 

    \begin{itemize}[label=$\circ$]
  	  \item Une humidité atmosphérique élevée peut réduire la demande évaporative à certaines périodes.
  	
  	  \item Une forte variabilité spatiale et temporelle des pluies.
  	  
      \item Un rayonnement élevé et des températures élevées pendant les périodes sèches.
  	  
      \item Une couverture nuageuse qui limite l’utilisation régulière des images satellitaires optiques.
	\end{itemize}

Les systèmes les plus pertinents combinent les mesures locales et le bilan hydrique pour éviter les erreurs d’irrigation liées aux pluies non enregistrées ou aux mesures non représentatives d’un capteur unique.
\end{itemize}

\begin{itemize}[label=$\bullet$]
   	\item \textbf{Travaux et orientations en zones sahéliennes}

En zone sahélienne, l’irrigation numérique répond à la rareté de l’eau, aux températures élevées et à la forte demande évaporative, notamment pour le maraîchage, les cultures de contre-saison et le pompage solaire. 
Les systèmes étudiés dans ces environnements cherchent généralement à :

    \begin{itemize}[label=$\circ$]
  	    \item Réduire les pertes d’eau dues aux apports excessifs.
  	  
  	    \item Ajuster les doses aux différents stades culturaux.

        \item Prévenir le stress hydrique pendant les périodes critiques.

        \item Adapter le fonctionnement des pompes solaires à la disponibilité énergétique.
    \end{itemize}

Les principaux défis sont l’ensoleillement intense, la poussière, la chaleur extrême, l’accès limité à l’électricité, au réseau et à la maintenance, ainsi que l’entretien des équipements, le vieillissement des batteries et la dérive des capteurs.
\end{itemize}

\begin{itemize}[label=$\bullet$]
   	\item \textbf{Enseignement des travaux antérieurs}

Les systèmes d’irrigation intelligents associent capteurs d’humidité, stations météo, calcul de l’évapotranspiration, communications sans fil et plateformes de traitement. 
Ils évoluent vers des approches combinant plusieurs variables et des modèles prédictifs. Les enseignements sont :  

    \begin{itemize}[label=$\circ$]
  	    \item La calibration locale selon la texture, la salinité et la densité apparente du sol.
  	
  	    \item La redondance : plusieurs capteurs et plusieurs profondeurs améliorent la fiabilité.
  	
        \item Validation agronomique : économie d’eau sans baisse excessive du rendement ou de la qualité.
    
        \item Mode hors ligne : maintien des calculs essentiels et du stockage en cas d’interruption du réseau.

        \item Le coût total doit être évalué : il faut considérer l’achat, l’installation, la connectivité, le remplacement des batteries, la calibration et la maintenance.
	\end{itemize}

\end{itemize}









   
 





    









 





\section{Salinité des sols et stratégies de correction}
\subsection{Problématique de la salinité des sols dans la zone d'étude}
La salinité des sols est l’un des principaux facteurs de dégradation des terres dans le Sine-Saloum, particulièrement dans la région de Fatick. Selon Amar \cite{noauthor_salinisation_nodate}, plus de 1,7 million d'hectares de terres arables au Sénégal sont affectées par le sel, conséquence de l’intrusion marine, de la remontée capillaire des sels et de  la variabilité climatique. Dans les zones côtières, 
la salinité transforme progressivement les rizières et les champs maraîchers en tannes improductives, réduisant ainsi la disponibilité des terres cultivables. Cette situation compromet la sécurité alimentaire et pousse de nombreux producteurs à abandonner leurs terres.


\subsection{Effet du stress salin sur les cultures maraîchères}
Le stress salin agit à plusieurs niveaux sur les cultures maraîchères. D’une part, il réduit la germination des semences et ralentit la croissance des jeunes plants. D’autre part, il provoque une toxicité ionique due à l'accumulation de sodium ($\mathrm{Na^{+}}$) et de chlorure ($\mathrm{Cl^{-}}$), qui perturbe l’absorption des nutriments essentiels comme le potassium et le calcium. Dans les Niayes et le bassin arachidier, les cultures de tomates, oignons et chou présentent des baisses de rendement pouvant dépasser 50\% \cite{fall_salt_nodate}. Les observations de Fall Ba et al. \cite{fall_salt_nodate} montrent que la salinité  peut entraîner une réduction significative de la photosynthèse et est associé à une augmentation du flétrissement prématuré, ce qui fragilise la rentabilité des exploitations maraîchères.

\subsection{Cultures étudiées et sensibilité à la salinité}
Trois cultures maraîchères ont été retenues : la laitue (Lactuca sativa), le navet (Brassica rapa) et le gombo (Abelmoschus esculentus). Le tableau ci-dessous présente les caractéristiques utiles à cette étude.

{\small
\begin{longtable}{|p{2cm}|p{3.2cm}|p{4.4cm}|p{4.4cm}|}
    \caption{Cycle, sensibilité au sel et besoins en eau des trois cultures}
    \label{tab:cultures_sel_eau} \\
    \hline
    \textbf{Culture} & \textbf{Cycle} & \textbf{Sensibilité au sel} & \textbf{Besoins en eau} \\
    \hline
    \endfirsthead
    \hline
    \textbf{Culture} & \textbf{Cycle} & \textbf{Sensibilité au sel} & \textbf{Besoins en eau} \\
    \hline
    \endhead
    Laitue & 40 à 50 jours du semis à la récolte & Sensibilité modérée, seuil de 1,3 à 1,75~dS/m ; effet principal : osmotique & Arrosage régulier et contrôlé (racines très limitées) ; manque : arrêt de croissance ; excès : asphyxie des racines \\
    \hline
    Navet & 45 à 48 jours du semis à la récolte & Moins sensible que la laitue mais moins résistant que le gombo & Irrigation régulière, surtout à l'expansion de la racine ; supporte mal l'alternance sécheresse et saturation \\
    \hline
    Gombo & 2 à 3 mois & Plus résistant que la laitue et le navet & Besoins constants, surtout à la floraison et à la fructification ; un déficit provoque la chute des fleurs et des fruits \\
    \hline
\end{longtable}
}

Ces exigences, en particulier la sensibilité au sel et la régularité des apports en eau, justifient le pilotage de l’irrigation étudié dans ce mémoire.

\subsection{Amendements organiques utilisés pour lutter contre la salinité (biochar, compost, coque d’arachide et biofertilisant)}
Face à ce problème, les amendements organiques apparaissent comme une solution durable et adaptée au contexte sénégalais. Le biochar, obtenu par  pyrolyse de biomasse, améliore la capacité de rétention d'eau et favorise la fixation des sels, réduisant ainsi leur disponibilité pour les plantes. Les coques d’arachides, abondantes dans le bassin arachidier, sont utilisées comme paillage ou compost. Elles contribuent à diminuer la conductivité électrique des sols et à améliorer leur structure. Le compost, riche en MO, stimule l'activité microbienne et augmente la tolérance des cultures au stress salin. Enfin, les fertilisants apportent des  micro-organismes bénéfiques qui favorisent la croissance des plantes en conditions difficiles. Les travaux de Diédhiou et al. \cite{diedhiou_investigation_2023} en Casamance ont montré que l’application combinée de biochar et de compost améliore la survie et le rendement du riz en condition saline, avec des taux de survie supérieurs à 90\%.

\subsection{Travaux antérieurs sur l'efficacité de ces amendements}
Les recherches menées par Niang et al. \cite{niang_degradation_2017} dans le Haut Saloum et le delta du Sénégal ont démontré que l’apport de compost et l’utilisation de CA permettent une réhabilitation progressive des sols salés. Ces pratiques agroécologiques favorisent la restauration de la fertilité et la reprise des cultures maraîchères. De plus, les initiatives locales de l’ISRA et du CIRAD dans le département de Fatick ont confirmé que la combinaison d’amendements organiques et de reboisement avec des espèces halophytes (Casuarina, Melaleuca) constitue une approche durable pour restaurer les sols et améliorer la résilience des systèmes agricoles. \textbf{Ces résultats montrent que les solutions endogènes, adaptées aux réalités locales, sont plus efficaces que les interventions purement chimiques}.



\section{Conclusion}
L’analyse des travaux antérieurs montre que la gestion durable de l’arrosage nécessite la combinaison des données météorologiques, de l’humidité  du sol et des caractéristiques des cultures. Les méthodes basées sur l'évapotranspiration, notamment FAO-56 Penman-Monteith \cite{allen_crop_1998-1} et le coefficient Kc, constituent des bases pour estimer les besoins en eau. Toutefois, leur adaptation aux conditions locales, ainsi que la fiabilité et le coût des capteurs, restent des défis importants. La salinité des sols constitue également une contrainte majeure nécessitant des stratégies de correction adaptées. Ces constats justifient ainsi la mise en place d’une solution numérique intégrée, adaptée au contexte local, pour améliorer la gestion de l’irrigation et favoriser une agriculture plus durable.
\chapter{PROTOCOLE EXPÉRIMENTAL}
\section{Introduction}
Ce chapitre présente le protocole expérimental mis en place pour comparer l'agriculture traditionnelle et l'agriculture de précision. Il s'organise autour de quatre volets : le dispositif expérimental mis en place, le matériel végétal, les paramètres du capteur et méthodologie de calcul, et enfin le processus d'aménagement des plots avant une conclusion partielle. 
L'objectif est de définir un cadre méthodologique rigoureux, permettant de mesurer et de comparer objectivement les performances des deux systèmes de gestion des ressources hydriques et nutritives.

\section{Dispositif expérimental}

\subsection{Présentation du site expérimental}
L'Université du Sine Saloum El Hadj Ibrahima Niass (USSEIN)  dispose de plusieurs campus, dont son siège principal situé à Kaolack, dans le quartier de Sing Sing, à proximité directe de l'école franco-arabe sur la route de Sing Sing. C'est dans ce cadre que le champ d'application Agro TIC a été mis en place vers 2020--2021, en tant qu'espace expérimental dédié aux activités d'enseignement, de recherche et de démonstration en agriculture numérique ; il se situe au Sud-Est de l'université, et c'est à l'intérieur de ce champ, au fond, que se trouve notre parcelle expérimentale, dont le schéma ci-dessous présente une représentation. 
\begin{figure}[H]
    \centering
    \imgstd{images/03_site_experimental.jpg}
    \caption{Représentation cartographique du site expérimental. }
    \label{fig:site2}
\end{figure}

\subsection{Représentation de la surface expérimentale}

Notre surface d'expérimentation est représentée par le schéma ci-dessous. Elle est composée de 6 plots, répartis en 3 plots pour la zone test et 3 plots pour la zone témoin, chaque zone comportant un plot de gombo, un plot de navet et un plot de laitue. Les espacements entre les cultures varient selon les espèces : 30 cm pour le gombo, et 15 cm aussi bien pour le navet que pour la laitue. Pour la collecte des données, un capteur 8-en-1, fournissant huit paramètres du sol, est branché manuellement sur chaque plot test trois fois par jour, tandis qu'un capteur DHT22 vient compléter ce dispositif en mesurant la température et l'humidité ambiante du milieu.
\begin{figure}[H]
    \centering
    \imgstd{images/04_surface_experimentation.jpg}
    \caption{Représentation de la surface d'expérimentation. }
    \label{fig:surface2}
\end{figure}

\subsection{Méthode d’irrigation et d’apport en nutriment} 
\textbf{Zone Test :}

Dans la zone test, l'irrigation et l'apport en nutriments reposent sur une mesure des besoins réels de chaque culture. Trois plots y sont aménagés : un pour la laitue, un pour le navet et un pour le gombo chacun équipé d'un capteur mesurant huit paramètres du sol (humidité, température, teneur en sel, conductivité, pH, N, P et K), associé a l'utilisation de AccuWeather. Ce capteur, connecté à la machine, permet des relevés manuels deux fois par jour, facilitant ainsi la lecture des données. Les informations recueillies sont intégrées à la méthode de Penman-Monteith \cite{allen_crop_1998-1} conformément à la méthode standardisée FAO-56 \cite[chap.~2]{allen_crop_1998-1} afin d'estimer les besoins en eau des plantations et de déterminer en temps réel une éventuelle insuffisance ou un surplus d'eau ou d'éléments nutritifs. Ce calcul sert ainsi d'aide à la décision : l'irrigation et l'apport en nutriments ne sont effectués manuellement qu'en cas de besoin réel de la plante, de façon mesurée. 

 \textbf{Zone témoin :}
 
Dans la zone témoin, l'irrigation suit une méthode traditionnelle fixe, indépendante des besoins réels des plantes. Trois plots de $2\,\mathrm{m^2}$ y sont également aménagés : un pour la laitue, un pour le navet et un pour le gombo. L'arrosage y est effectué manuellement, chaque jour, matin et soir, selon un calendrier fixe qui ne tient compte ni des caractéristiques du sol ni de celles du climat, ni des besoins spécifiques de chaque culture.




\section{Matériel végétal}
Trois cultures ont été étudiées : 

\begin{table}[H]
    \centering
    \caption{Variétés et modes de semis des trois cultures}
    \label{tab:varietes_semis}
    \begin{tabular}{|p{3cm}|p{5cm}|p{5.5cm}|}
        \hline
        \textbf{Culture} & \textbf{Variété} & \textbf{Mode de semis} \\
        \hline
        Laitue & EDEN & repiquage en motte \\
        \hline
        Navet & SUPER LONGO 2000 & Semis direct \\
        \hline
        Gombo & CLEMSON SPINELESS & Semis direct \\
        \hline
    \end{tabular}
\end{table}

\begin{figure}[H]
    \centering
    \imgstd{images/05_pepiniere_laitue.jpg}
    \caption{Pépinière de laitue variété EDEN. }
    \label{fig:laitue2}
\end{figure}

\begin{figure}[H]
    \centering
    \imgstd{images/06_semence_navet.jpg}
    \caption{Semence de Navet SUPER LONGO 2000. }
    \label{fig:navet2}
\end{figure}

\begin{figure}[H]
    \centering
    \imgstd{images/07_semence_gombo.jpg}
    \caption{Semence de Gombo CLEMSON SPINELESS. }
    \label{fig:gombo2}
\end{figure}

\subsection{Aménagement du sol et amendements}
Les amendements n'ont pas été pesés avec précision. Pour l'ensemble des trois plots TS, les quantités utilisées ont été les suivantes : 
un sac de compost conditionné dans un sac de riz de 50 kg, un demi-sac de biochar et de PCA, dont seule la moitié a été utilisée ; et 2 à 3 kg d'EGB. 
Le compost a été apporté en deux temps : une partie lors de l'aménagement des plots et une autre au cours de l'expérimentation.


\section{Paramètres capteurs et méthodologie de calcul}
\subsection{Paramètres capteurs}
La collecte des données édaphiques a été réalisée à l’aide d’un capteur 8-en-1 mesurant la température, l’humidité, le $pH$, la $EC$ et les teneurs en $N$, $P$ et $K$. 
Les données climatiques, notamment la température, l’humidité de l’air et la vitesse du vent, ont été relevées à partir d’une application météorologique (AccuWeather).

Les paramètres mesurés ont servi au suivi des conditions du sol, tandis que les données climatiques ont été utilisées pour estimer les besoins en eau des cultures.

\begin{longtable}{|p{0.8cm}|p{3cm}|p{3cm}|p{1.8cm}|p{5cm}|}
    \caption{Capteurs et paramètres mesurés sur la parcelle expérimentale}
    \label{tab:capteurs_parametres} \\
    \hline
    \textbf{N°} & \textbf{Capteur / Instrument} & \textbf{Paramètre mesuré} & \textbf{Unité} & \textbf{Utilisation dans les calculs} \\
    \hline
    \endfirsthead
    \hline
    \textbf{N°} & \textbf{Capteur / Instrument} & \textbf{Paramètre mesuré} & \textbf{Unité} & \textbf{Utilisation dans les calculs} \\
    \hline
    \endhead
    1 & \multirow{3}{3cm}{Application météo} & Température de l'air & °C & Calcul de $ET_0$ (formule de Penman-Monteith) \\
    \cline{1-1}\cline{3-5}
    2 & & Humidité relative de l'air & \% & Calcul de $e_s$, $e_a$ et du déficit de pression de vapeur \\
    \cline{1-1}\cline{3-5}
    3 & & Vitesse du vent & m/s & Calcul de $ET_0$ et de l'évapotranspiration potentielle \\
    \hline
    4 & \multirow{7}{3cm}{Capteur 8-en-1} & Température du sol & °C & Calcul du flux thermique du sol \\
    \cline{1-1}\cline{3-5}
    5 & & Humidité volumique du sol & \% (m\textsuperscript{3}/m\textsuperscript{3}) & Suivi du déficit hydrique, déclenchement des irrigations \\
    \cline{1-1}\cline{3-5}
    6 & & pH du sol & --- & Évaluation de la disponibilité des nutriments, ajustement de la fertilisation \\
    \cline{1-1}\cline{3-5}
    7 & & Conductivité électrique du sol & \textmu S\,cm$^{-1}$ & Contrôle de la salinité, adaptation des doses de fertilisation \\
    \cline{1-1}\cline{3-5}
    8 & & Teneur en azote (N) & mg/kg & Diagnostic de fertilité, correction et ajustement des apports d'intrant \\
    \cline{1-1}\cline{3-5}
    9 & & Teneur en phosphore (P) & mg/kg & Diagnostic de fertilité, correction et ajustement des apports d'intrant \\
    \cline{1-1}\cline{3-5}
    10 & & Teneur en potassium (K) & mg/kg & Diagnostic de fertilité, correction et ajustement des apports d'intrant \\
    \hline
    11 & Calculé via une formule & Ensoleillement ou rayonnement solaire & W/m\textsuperscript{2} & Calcul du rayonnement net ($R_n$) et de $ET_0$ \\
    \hline
\end{longtable}

\subsection{Besoins des cultures et Méthodologie de calcul}
\subsubsection{Besoins en eau}
Les besoins en eau des cultures ont été estimés à partir de l’$ET_0$ et du $K_c$, selon la relation présentée au chapitre 1 : 

\begin{equation}
    ET_c = K_c \times ET_0
    \label{eq:etc2}
\end{equation}

L’$ET_0$ a été calculée selon la méthode FAO-56 Penman-Monteith \cite{allen_crop_1998-1}, à partir des données météorologiques recueillies dans AccuWeather et des variables intermédiaires nécessaires au calcul. Les calculs ont été réalisés dans Excel.


\subsubsection{Variables climatiques utilisées}
Les données climatiques utilisées comprenaient les températures maximale et minimale, l’humidité relative maximale et minimale et la vitesse du vent. 
Le rayonnement solaire a été estimé à l’aide de la formule de Hargreaves. Ces données ont servi au calcul journalier de $ET_0$, puis à l’estimation de $ET_c$ pour chaque spéculation.


\subsubsection{Coefficients culturaux et stades de développement}
Les $K_c$ ont été appliqués en fonction des cultures étudiées et de leurs stades de développement. Les valeurs retenues pour le gombo, la laitue et le navet sont présentées dans le tableau ci-dessous : 

\begin{table}[H]
    \centering
    \caption{Coefficients culturaux ($K_c$) par culture et par stade de développement}
    \label{tab:coefficients_kc}
    \begin{tabular}{|p{3cm}|p{3cm}|p{4cm}|p{3cm}|}
        \hline
        \textbf{Culture} & \textbf{$K_c$ initial} & \textbf{$K_c$ milieu de saison} & \textbf{$K_c$ final} \\
        \hline
        Laitue & 0,70 & 1,05 & 0,95 \\
        \hline
        Navet & 0,70 & 1,00 & 0,95 \\
        \hline
        Gombo & 0,40 & 1,00 & 0,75 \\
        \hline
    \end{tabular}
\end{table}

Les stades de développement et leurs durées indicatives ont servi à déterminer le $K_c$ applicable à chaque période du cycle. 
Les besoins en eau ainsi estimés ont ensuite été utilisés pour déterminer les doses d’irrigation.

\subsection{Besoin en élément nutritif}
Ces cultures maraîchères ont des exigences nutritives spécifiques. Pour maximiser leur productivité, il est essentiel de trouver un équilibre entre les nutriments en $N$, en $P$, et en $K$. 
Pour ce faire, nous avons opté systématiquement pour une fertilisation de fond utilisant de la matière organique bien décomposée et du compost liquide à base de fientes de volaille. Toutefois, chaque culture requiert un ou des besoins nutritifs dominants, le tableau suivant donne les valeurs ($N$, $P$, $K$) seuils déclenchant la fertilisation pour chaque culture :

\begin{table}[H]
    \centering
    \caption{Valeurs seuils déclenchant la fertilisation sur les plots TS}
    \label{tab:seuils_fertilisation_ts}
    \begin{tabular}{|p{3cm}|p{3.4cm}|p{3.8cm}|p{3.8cm}|}
        \hline
        \textbf{Culture} & \textbf{Azote (N) (mg/kg)} & \textbf{Phosphore (P) (mg/kg)} & \textbf{Potassium (K) (mg/kg)} \\
        \hline
        Gombo & 40 & 32,5 & 125 \\
        \hline
        Navet & 40 & 32,5 & 125 \\
        \hline
        Laitue & 60 & 32,5 & 125 \\
        \hline
    \end{tabular}
\end{table}

\section{Processus d'aménagement des plots : étapes, rôle et justification des apports}
\subsection{Les étapes de l'aménagement}
\subsubsection{Désherbage et Isolement des herbes}
\begin{itemize}[label=$\bullet$]
   	 \item Le désherbage consiste à nettoyer les herbes sur une parcelle avant la préparation du sol. Cette action permet d’assainir la surface, facilitant les travaux de labour.
	\end{itemize}

\begin{itemize}[label=$\bullet$]
   	 \item L’isolement des herbes consiste à regrouper et retirer les herbes  hors de la zone de travail afin de maintenir la zone propre.
	\end{itemize}

    \begin{figure}[H]
    \centering
    \imgstd{images/08_desherbage_isolement_herbes.jpg}
    \caption{Processus de désherbage et d’isolement des herbes : début et fin. }
    \label{fig:desherbage2}
\end{figure}

\subsubsection{Le labour et la délimitation des plots}
Après un labour profond sur les trois plots constituant la zone test, nous avons ajouté du biochar, de la PCA, de l'EGB et du compost que nous avons soigneusement mélangés avant d'aplanir et de fixer les ados. En revanche, dans la zone témoin, nous avons simplement délimité et fixé les ados après un labour approfondi.

\begin{itemize}[label=$\bullet$]
\item Le labour est une pratique agricole qui consiste à retourner la couche superficielle de la terre. Cette opération permet d'aérer et d’ameublir le sol pour préparer les cultures à venir.

\item Le travail en profondeur mené lors de la réalisation des plots a pour but d'ameublir le sol, d'améliorer son aération et sa teneur en eau, tout en stimulant la croissance des racines des cultures à venir et l'intégration des amendements ajoutés.

\item La confection de plots avec des ados a contribué à délimiter les espaces cultivés et les voies de passage, rendant par conséquent les opérations et la supervision de l'irrigation plus aisées.
\end{itemize}

\begin{figure}[H]
    \centering
\imgstd{images/11_mise_en_place_plot_irrigation_fixer_ados.jpg}
    \caption{Mise en place des plots et irrigation pour la fixation des ados. }
    \label{fig:plotAdos2}
\end{figure}

\subsection{Rôle et justification du choix des apports(mécanisme attendu)}
Les éléments apportés sont généralement des éléments de correction du sol:
\subsubsection{Le biochar}
Le biochar possède une surface poreuse et aromatique riche en groupes fonctionnels chargés 
négativement (carboxyles \ce{-COOH}, hydroxyles \ce{-OH}, phénols \ce{Ar-OH}), ce qui lui donne 
une capacité d'échange cationique (CEC) élevée \cite{chandrasekaran_comprehensive_2026}. Cette CEC lui permet d'attirer
et de retenir des ions chargés positivement appelés cations du sol, dont le \ce{Na+} (sodium), qui est
principalement responsable de la salinité. Le mécanisme précis pour lutter contre la salinité 
repose sur un échange préférentiel : les cations polyvalents présents dans le biochar, notamment 
le calcium (\ce{Ca^2+}) et le magnésium (\ce{Mg^2+}), peuvent remplacer le \ce{Na+} sur les sites 
d'échange, permettant ensuite à l'excès de sodium libéré d'être facilement lessivé grâce à 
l'irrigation ou au ruissellement \cite{yuan_biochar_2023}. De plus, le biochar favorise l'absorption du potassium par la plante
tout en diminuant celle du sodium, participant à l'équilibre ionique et à la régulation osmotique de la plante \cite{chandrasekaran_comprehensive_2026,zubairu_overview_2025}.

    \begin{figure}[H]
    \centering
\imgstd{images/12_biochar.jpg}
    \caption{Poudre de charbon de bois (dit Biochar en agriculture). }
    \label{fig:biochar2}
\end{figure}

\subsubsection{La poudre de coque d’arachide}
La CA est principalement composée de 45~\% de cellulose, de 6~\% d’hémicellulose et de
36~\% de lignine, ce qui en fait un bon amendement pour le sol; sa teneur élevée en lignine
entraîne cependant une décomposition lente dans l’environnement naturel \cite{sawargaonkar_valorization_2024}. Grâce à sa
décomposition lente, elle représente une source durable de matière organique qui améliore
\textbf{l’aération et la structure} du sol sur le long terme, contrairement à un apport nutritif rapide
comme le compost. Une étude menée spécifiquement sur des sols salins a montré que
l’amendement avec des CA (en quantités de 4, 6 ou 8 tonnes par hectare) améliore les
propriétés du sol et la croissance de jeunes plants d’espèces ligneuses cultivées sur des
sols salins \cite{fall_effect_2018}. 

NB : \textbf{Les coquilles d’arachide représentent 30~\% du poids des arachides.}

\begin{figure}[H]
    \centering
\imgstd{images/13_coque_arrachide.jpg}
    \caption{Poudre de coque d’arachide. }
    \label{fig:arrachide2}
\end{figure}

\subsubsection{Le compost}
Le compost améliore les propriétés physiques, chimiques et biologiques du sol : il augmente la capacité de rétention en eau du sol, 
la disponibilité des nutriments nécessaires ($N$, $P$, $K$) et renforce la résistance des cultures face aux maladies \cite{goldan_assessment_2023}. 
Contrairement à un engrais minéral à libération immédiate, le compost libère ses éléments nutritifs de façon progressive, assurant ainsi un apport constant au sol et réduisant le besoin de fertilisations répétées.

En ce qui concerne plus spécifiquement la salinité, une étude portant sur l’irrigation
avec une eau saumâtre a montré qu’un compost enrichi biologiquement a contribué à
la réduction de l’accumulation de sodium de 60~\% et celle du chlore de 70~\%, tout en
améliorant la rétention d’eau et la structure du sol \cite{suvendran_soil_2025}.

\begin{figure}[H]
    \centering
\imgstd{images/14_compost.jpg}
    \caption{Compost (Matière organique). }
    \label{fig:compost2}
\end{figure}

\subsubsection{Le biofertilisant anti-sel ou engrais granulé bio}
Contrairement aux trois amendements précédents, dont l’action est essentiellement physicochimique, le EGB agit biologiquement, via des micro-organismes rhizosphériques. 
Les rhizobactéries halotolérantes favorisant la croissance des plantes (HT-PGPR, Halotolerant Plant Growth-Promoting Rhizobacteria) agissent par plusieurs mécanismes combinés : 
la production de métabolites secondaires et de sidérophores, la synthèse de phytohormones et d’enzymes, le maintien de l’équilibre ionique, et l’amélioration de la disponibilité en nutriments \cite{ramasamy_coping_2023}.

Le mécanisme d’équilibrage ionique est particulièrement efficace pour la lutte contre la
salinité : sous stress salin, les plantes réduisent leur évapotranspiration en renforçant l’absorption 
de potassium (K+) et en diminuant celle du sodium (Na+), ce qui atténue le stress osmotique et ionique. 
Ce processus est facilité par les bactéries PGPR qui augmentent l’absorption des nutriments bénéfiques par la plante \cite{gao_progress_2022}.

\begin{figure}[H]
    \centering
\imgstd{images/15_anti_sel.jpg}
    \caption{Engrais granulé bio. }
    \label{fig:EGB2}
\end{figure}
    

    

  




 





 


\section{Conclusion}
Ce chapitre a défini le cadre méthodologique de l'étude comparative entre l'agriculture de précision et l'agriculture traditionnelle. Le dispositif repose sur deux zones de trois plots de 2~m\textsuperscript{2} (gombo, navet, laitue) : une zone test, pilotée à l'aide du capteur 8-en-1 et du DHT22, et une zone témoin, irriguée selon un calendrier fixe. Les besoins en eau sont estimés par la méthode FAO-56 Penman-Monteith ($ET_c = ET_0 \times K_c$) et les besoins en nutriments par le calcul du déficit en N, P et K. Enfin, la zone test a bénéficié d'un amendement du sol (biochar, poudre de coque d'arachide, compost, engrais granulé bio) pour lutter contre la salinité. Ce protocole constitue la base de la mise en œuvre sur le terrain, présentée au chapitre suivant.

\chapter{MISE EN ŒUVRE RÉELLE ET SUIVI DE L'EXPÉRIMENTATION}
\section{Introduction}
Après avoir présenté le cadre théorique et méthodologique de cette étude, ce chapitre est consacré à la mise en œuvre de l’expérimentation ainsi qu’à son suivi tout au long du cycle cultural. Il décrit les différentes étapes de la conduite de l’essai, depuis l’aménagement des parcelles et l’installation des cultures jusqu’au suivi de leur évolution. Une attention particulière est accordée au dispositif d’irrigation, aux outils et méthodes de collecte des données, ainsi qu’à la gestion hebdomadaire des parcelles. Les approches déployées pour assurer la protection des cultures contre les maladies et les ravageurs sont également présentées. Finalement, l’évolution des divers plots au cours de l’expérimentation est illustrée afin de mettre en évidence les circonstances dans lesquelles les données ont été recueillies. L'ensemble de ces éléments constitue une base essentielle à la compréhension et à l'interprétation des résultats présentés dans le chapitre suivant.
\section{Mise en place des plots et des cultures}
\subsection{Mise en place des plots}
La mise en place des plots renvoie à la délimitation et au nivellement des 6 plots après un labour profond. Pour délimiter nos plots nous avons utilisé la méthode 3-4-5. C'est est une méthode simple basée sur le théorème de Pythagore pour tracer des angles droits (90 degrés). En agriculture, elle est très utile pour délimiter correctement les planches, les parcelles et les lignes de culture. Dans notre cas, pour l'obtention de tous les sommets, nous avons utilisé un piquet pour marquer un point A, puis avec le ruban-mètre nous avons mesuré 3~m pour avoir le point B. À partir de A, dans une direction perpendiculaire, nous mesurons 4~m pour y placer le piquet C, puis à partir de ces 3 points nous projetons avec des arcs de cercle afin d'obtenir les 4 points déterminant ainsi la surface totale de nos plots expérimentaux.
Après la délimitation des plots, nous avons élevé les ados, puis nous avons bien nivelé et enfin nous avons fixé les ados en les arrosant. Ce schéma ci-dessous est une parfaite illustration de la méthode 3-4-5.

\begin{figure}[H]
    \centering
    \imgstd{images/16_methode_3_4_5.png}
    \caption{Illustration de la méthode 3-4-5 pour la délimitation de la surface des plots. }
    \label{fig:plot3-4-5}
\end{figure}

La figure ci-dessous présente la configuration obtenue après utilisation de la méthode 3-4-5 pour la délimitation de notre surface.

\begin{figure}[H]
    \centering
    \imgstd{images/17_fin_confection.jpg}
    \caption{Fin de confection des 6 plots. }
    \label{fig:confection3}
\end{figure}

\subsection{Mise en place des cultures sur les plots}

\begin{itemize}[label=$\bullet$]
   	 \item \textbf{Organisation des cultures sur les plots}

     Nous avons 6 plots constitués de 3 plots témoins et de 3 plots tests. Par annotation, de la gauche vers la droite, nous avons sur la première ligne les plots de gombo, au milieu les plots de navet et à droite les plots de laitue. Chaque plot test est aligné sur la même ligne avec son témoin de la même spéculation. Le plot gombo test est nommé GOM TS et son témoin GOM T, le plot test navet est nommé NAV TS et son témoin NAV T, et enfin, le plot laitue test est nommé LAI TS et son témoin LAI T.
	\end{itemize}
\begin{itemize}[label=$\bullet$]
   	 \item \textbf{Mise en place des cultures}

     Nous avons réalisé un semis direct de gombo sur les plots GOM TS et GOM T, et de navet sur les plots NAV TS et NAV T, le soir du 25 juin 2026. Pour les plots de laitue, un repiquage a été effectué le soir du 25 juillet 2026.
	\end{itemize}

\begin{itemize}[label=$\bullet$]
   	 \item \textbf{Méthodes de semis et de repiquage}

     Pour le gombo, nous avons semé par poquets. Dans chaque poquet, nous avons mis une ou deux graines, avec une profondeur égale à deux fois la taille de la graine. Le semis du navet est de type semis en ligne, avec une ou deux graines par trou et une distanciation de 15 cm. La laitue a également été repiquée en ligne, avec espacement de 15 cm.
	\end{itemize}






\section{Irrigation}
Pour piloter l’irrigation, nous procédons par détermination de $ET_c$, qui dépend de $ET_0$ du jour et du $K_c$ de chaque culture. Les $K_c$ de nos trois cultures sont déjà définis dans notre protocole, ils se diffèrent selon la spéculation et le stade de développement des cultures. Au début de notre expérimentation, nous faisions le calcul de $ET_c$ manuellement, puis nous avons établi un classeur Excel dans lequel les calculs sont effectués automatiquement dès que les paramètres météorologiques $T_{max}$, $T_{min}$, $RH_{max}$, $RH_{min}$ et la vitesse du vent $U_2$ sont renseignés, ainsi que la date du jour correspondante. La date nous renseigne sur le numéro du jour de l’année correspondant. Cela nous aide à déterminer le rayonnement solaire à partir de la formule de Hargreaves \cite{koc_reference_2023}.

\subsection{Détermination manuelle des paramètres pour l’irrigation}
Ici nous avons donné les formules principales et les détails nous les avons en annexe (voir l’Annexe~\ref{ann:et0}).
\subsubsection{Détermination de \texorpdfstring{$ET_0$}{ET0}}
L’évapotranspiration de référence $ET_0$ est calculée selon la méthode de Penman-Monteith recommandée par la FAO-56 \cite[chap.~3]{allen_crop_1998-1}

\begin{equation}
    ET_0 = \frac{0{,}408 \times \Delta \times (R_n - G) + \gamma \times \dfrac{900}{T + 273} \times u_2 \times (e_s - e_a)}{\Delta + \gamma \times (1 + 0{,}34 \times u_2)}
    \label{eq:et0_final}
\end{equation}

Où :
\begin{itemize}[label=$\bullet$]
   	\item $ET_0$ : évapotranspiration de référence (mm/jour) ;
 
    \item $\Delta$ : pente de la courbe de pression de vapeur saturante (kPa/°C) ;

    \item $R_n$ : rayonnement net (MJ/m²/jour) ;

    \item $G$ : flux de chaleur dans le sol (MJ/m²/jour), considéré comme négligeable à l’échelle journalière ;

    \item $\gamma$ : constante psychrométrique (kPa/°C) ;

    \item $T$ : température moyenne journalière (°C) ;

    \item $u_2$ : vitesse du vent à 2 m de hauteur (m/s) ;

    \item $e_s-e_a$ : déficit de pression de vapeur (kPa).

\end{itemize}
Le calcul nécessite les données météorologiques suivantes : $T_{max}$ et $T_{min}$, $RH_{max}$ et $RH_{min}$ et vitesse du vent. La latitude et l’altitude du site sont aussi utilisées dans les calculs des paramètres intermédiaires.

Le rayonnement extraterrestre $R_a$ est déterminé à partir de la latitude et du jour de l’année, tandis que le rayonnement solaire $R_s$ est estimé par la formule de Hargreaves \cite{koc_reference_2023}, en utilisant un coefficient $K_{Rs}=0,16$.
Le rayonnement net $R_n$, les pressions de vapeur $e_s$ et $e_a$, le déficit de pression de vapeur et les autres paramètres sont ensuite déterminés conformément à la méthode FAO-56 \cite[chap.~3]{allen_crop_1998-1}

\subsubsection{Détermination de \texorpdfstring{$ET_c$}{ETc} en mm/jour}
L’évapotranspiration culturale est calculée à partir de l’$ET_0$ et du coefficient cultural $K_c$ :

\begin{equation}
    ET_c = ET_0 \times K_c \quad \text{mm/jour}
    \label{eq:etc_final}
\end{equation}
Où :
\begin{itemize}[label=$\bullet$]
   	\item $ET_c$ : évapotranspiration culturale (mm/jour) ;

    \item $ET_0$ : évapotranspiration de référence (mm/jour) ;

    \item $K_c$ : coefficient cultural correspondant à la culture et à son stade de développement.
\end{itemize}

\subsubsection{Détermination de la dose théorique d’irrigation}
Pour connaître la dose d'irrigation journaliere en litres, nous avons multiplié $ET_c$ par la surface de chaque plot. Pour notre expérimentation avec des plots de 2~m\textsuperscript{2}, nous avons :

\begin{equation}
    D_i = ET_c \times 2
    \label{eq:dose_irrigation}
\end{equation}

Avec :
\begin{itemize}[label=$\bullet$]
   	\item $D_i$ : dose théorique journalière (L/jour) ;

    \item $ET_c$ : évapotranspiration culturale (mm/jour) ;

\end{itemize}

Cette conversion repose sur l’équivalence entre 1 mm d’eau et 1 L/m².

\subsection{Détermination automatique des paramètres de \texorpdfstring{$ET_0$}{ET0} sur Excel}
\subsubsection{Détermination automatique des paramètres sur Excel}
Pour éviter les erreurs de calcul et optimiser la gestion du temps, nous avons établi un classeur Excel comportant quatre feuilles permettant d'automatiser l'ensemble des calculs. 
Dans ce classeur, la latitude et l'altitude du site sont renseignées une seule fois, tandis que les paramètres indiqués en bleu, notamment la date, $T_{\max}$, $T_{\min}$, $RH_{\max}$, $RH_{\min}$ et $u_2$, sont renseignés quotidiennement.
Ici nous avons illustré avec une seule capture l'apercu du classeur et les détails sont présentés en annexe (voir l’Annexe~\ref{ann:et0}).

\begin{itemize}[label=$\bullet$]
   	\item \textbf{La feuille ET0 FAO-56 :}
\end{itemize}
    
    \begin{figure}[H]
    \centering
    \imgstd{images/18_ET0_partie1.jpg}
    \caption{Feuille de calcul automatique de $ET_0$ première partie. }
    \label{fig:partie1}
\end{figure}

\subsection{Détermination des paramètres hydriques du sol (\texorpdfstring{$H_{cc}$, $RFU$, $P$, $RU$}{Hcc, RFU, P, RU} et \texorpdfstring{$S_c$}{Sc})}
L’humidité à la capacité au champ ($H_{cc}$) correspond à la quantité d’eau retenue par un sol après saturation et après écoulement de l’eau excédentaire. 
Pour déterminer la $Hcc$ de notre sol, nous avons préalablement inondé les plots, puis nous les avons laissés pendant 24 heures. Après cette période, nous avons mesuré l’humidité du sol. \textbf{La valeur obtenue est de 28~\% (soit 0,28)}.

Le point de flétrissement permanent $H_{pf}$ correspond à la quantité minimale d’eau du sol en dessous de laquelle les plantes ne peuvent plus extraire suffisamment d’eau et subissent un flétrissement irréversible. Dans notre expérimentation, il a été estimé par une méthode de classe texturale, en utilisant des valeurs de référence issues des tables FAO ou USDA. Étant donné que notre sol est de texture \textbf{sablo-limoneux}, $H_{pf}$ =0,11 \cite[chap.~7,Table.~19]{allen_crop_1998-1}.

\subsubsection{Détermination de la réserve utile}
La réserve utile $RU$ représente la quantité d’eau que le sol peut retenir et restituer aux racines. Elle est calculée selon la formule suivante \cite[chap.~8]{allen_crop_1998-1}
La réserve facilement utilisable ($RFU$) représente la fraction de la $RU$ que la plante peut extraire du sol avant l’apparition d’un stress hydrique. 
La fraction d’épuisement $p$ correspond à la proportion de la $RU$ mobilisable avant l’apparition d’un stress hydrique. Sa valeur dépend notamment de la culture et de l’$ET_c$. 
Pour tenir compte des variations journalières de l’$ET_c$ au cours de l’expérimentation, la valeur de référence est ajustée selon la relation proposée par la FAO-56 \cite{allen_crop_1998-1}.

Le \textbf{seuil critique} $S_c$ correspond au niveau d’humidité du sol à partir duquel l’eau facilement disponible pour la culture devient insuffisante. Il constitue le seuil de déclenchement de l’irrigation afin d’éviter un stress hydrique. 
Le \textbf{stock actuel d’eau} $S_a$ fait référence à la quantité d’eau réellement présente dans la zone racinaire à un moment précis. Il représente l’indicateur principal du déclenchement ou non de l’irrigation.

{\small
\begin{longtable}{|p{3.6cm}|p{6.2cm}|p{4.2cm}|}
    \caption{Paramètres hydriques du sol et équations utilisées}
    \label{tab:parametres_hydriques} \\
    \hline
    \textbf{Paramètre} & \textbf{Méthode de détermination ou équation} & \textbf{Valeur ou précision} \\
    \hline
    \endfirsthead
    \hline
    \textbf{Paramètre} & \textbf{Méthode de détermination ou équation} & \textbf{Valeur ou précision} \\
    \hline
    \endhead
    Humidité à la capacité au champ ($H_{cc}$) & Inondation des plots, attente de 24 heures, puis mesure de l'humidité du sol & 0,28 (28~\%) \\
    \hline
    Point de flétrissement permanent ($H_{pf}$) & Estimation à partir de la classe texturale et des valeurs de référence FAO/USDA \cite[chap.~7, table~19]{allen_crop_1998-1} & 0,11 \\
    \hline
    Réserve utile ($RU$) & $RU = (H_{cc} - H_{pf}) \times Z_r \times 1000$ \cite[chap.~8]{allen_crop_1998-1} & Calculée selon la culture \\
    \hline
    Réserve facilement utilisable ($RFU$) & $RFU = p \times RU$ \cite[chap.~8]{allen_crop_1998-1} & Calculée selon la culture et $ET_c$ \\
    \hline
    Fraction de déplétion ajustée ($p$) & $p = p_{\text{tab}} + 0{,}04 \times (5 - ET_c)$ \cite[chap.~8, table~22]{allen_crop_1998-1} & Résultat limité entre 0,1 et 0,8 \\
    \hline
    Seuil critique ($S_c$) & $S_c = (H_{cc} \times Z_r \times 1000) - RFU$ & Seuil de déclenchement de l'irrigation \\
    \hline
    Stock actuel d'eau ($S_a$) & $S_a = \theta_{\text{actuel}} \times Z_r \times 1000$ & Exprimé en mm \\
    \hline
\end{longtable}
}

Avec : $Hcc$ : humidité du sol à la capacité au champ ; $H_{pf}$ : humidité du sol au point de flétrissement permanent ; $Z_r$ : profondeur racinaire, exprimée en mètres ; $p$ : fraction d’épuisement ajustée ; 
$p_{tab}$ : valeur de référence de la fraction d’épuisement selon la culture ; $ET_c$ : évapotranspiration culturale journalière, exprimée en mm/jour ; $\theta_{\text{actuel}}$ : l'humidité volumique du sol mesurée par le capteur 8-en-1 à l'instant t (en m³/m³).

Les profondeurs racinaires ont été déterminées à partir des valeurs de référence de la FAO-56, en fonction des cultures \cite[chap.~8,Table.~22]{allen_crop_1998-1}

\begin{table}[H]
    \centering
    \caption{Profondeurs racinaires de référence des trois cultures}
    \label{tab:profondeurs_racinaires}
    \begin{tabular}{|p{3cm}|p{3cm}|p{3cm}|p{3cm}|}
        \hline
        \textbf{Cultures} & \textbf{Laitue} & \textbf{Navet} & \textbf{Gombo} \\
        \hline
        $Z_r$ (m) & 0,35 & 0,50 & 0,60 \\
        \hline
    \end{tabular}
\end{table}

Les paramètres déterminés ont permis le suivi de l’état hydrique du sol et le pilotage de l’irrigation. 
La section suivante présente les principaux outils mobilisés pour l’acquisition des données météorologiques et agronomiques nécessaires au suivi de notre expérimentation.

\section{Présentation des outils d'acquisition de données}
Dans le cadre de notre étude, les outils mobilisés relèvent principalement des plateformes et des systèmes de collecte de données, notamment Google Earth, 
l’application météorologique du smartphone et les capteurs. Ces outils permettent d’obtenir des données exploitables nécessaires pour réaliser les divers calculs.

\subsection{Acquisition des données géographiques et météorologiques}
Les coordonnées géographiques (\textbf{latitude = 14 et altitude = 3}) du site expérimental ont été déterminées à l’aide de Google Earth.

Les données météorologiques nécessaires au calcul de $ET_0$ ($T_{max}$, $T_{min}$, $RH_{max}$, $RH_{min}$ et $U_2$) ont été relevées à partir de l’application AccuWeather.

Pour le développement de l’application du pilotage de l’irrigation et de la fertilisation, \textbf{l’acquisition des données météorologiques sera automatisée à l’aide de l’API OpenWeather}. 
Les méthodes de calcul de l’$ET_0$ et de l’$ET_c$ seront conservées ; seule la source des données météorologiques changera.

\subsection{Le capteur 8-en-1}
Le capteur 8-en-1 a été utilisé pour mesurer les paramètres du sol dans les différents plots, notamment la température, l’humidité, l'$EC$, le $pH$ et les teneurs en $N$, $P$ $K$. 
L’humidité a particulièrement contribué au suivi de l’état hydrique du sol, tandis que les autres paramètres ont permis de suivre son évolution physicochimique.

\begin{figure}[H]
    \centering
    \imgstd{images/26_mon_8_en_1.jpg}
    \caption{Capteur de sol 8-en-1. }
    \label{fig:notreCapteur}
\end{figure}
Les données météorologiques et les mesures du sol constituent les principales données utilisées pour le suivi de l’expérimentation et l’évaluation des besoins en eau des cultures.

\subsection{Caractérisation des paramètres du sol étudiés}
Les paramètres réellement étudiés sont : la température,l'humidité, le $pH$, le $EC$, $N$, $P$ et $K$. Ils sont mesurés à l’aide du capteur 8-en-1.

\subsubsection{La température du sol et l'humidité du sol}
La mesure de la température a permis de suivre son évolution au cours du cycle cultural. Les données collectées ont permis de caractériser les conditions du sol et de les mettre en relation avec les autres paramètres étudiés.

L’humidité du sol a permis de suivre l’évolution de la disponibilité en eau dans les plots tests. Les mesures ont permis d’apprécier les conditions hydriques du sol et d’adapter l’irrigation en fonction du $S_c$ défini, afin de maintenir des conditions favorables au développement des cultures.

\subsubsection{ Le  potentiel hydrogène du sol}
La mesure du $pH$ a permis de suivre son évolution au cours du cycle cultural. Les données collectées ont permis de caractériser les conditions physicochimiques du sol dans les différents plots et d’observer les variations éventuelles entre les plots T et les plots TS.

\subsubsection{L'électroconductivité ou conductivité électrique}
La mesure de la conductivité électrique (EC) a permis de suivre l’évolution des conditions physicochimiques du sol. Ces mesures ont complété les données relatives à l’humidité, à la température et au $pH$, notamment pour observer les variations de la teneur en sels solubles.

\subsubsection{L'azote, le phosphore et le potassium}
Le suivi de l’azote $N$, du phosphore $P$ et du potassium $K$ a permis d’obtenir des informations complémentaires sur l’état nutritif du sol et sur les conditions favorables au développement des cultures.

\section{Gestion hebdomadaire et suivi de l'expérimentation culturale}
La réussite de l’expérimentation repose sur un suivi régulier des cultures tout au long de leur cycle de développement. Ce suivi a porté sur l’irrigation, l’entretien des cultures, le suivi phytosanitaire et l’évolution des différents plots.

\subsection{Protocole de suivi des cultures}
Dès le lendemain du semis, le suivi des cultures a été effectué trois fois par jour (7~h–8~h, vers 13~h et 18~h–19~h). Pendant la première semaine, une irrigation régulière a été assurée jusqu’à la levée et à l’installation des jeunes plants. Les observations réalisées ont été renseignées dans l’agenda de culture et accompagnées de photographies.

\begin{table}[H]
    \centering
    \small
    \caption{Protocole de suivi des cultures durant l'expérimentation}
    \label{tab:protocole_suivi}
    \begin{tabular}{|p{3.2cm}|p{5.4cm}|p{5.4cm}|}
        \hline
        \textbf{Éléments suivis} & \textbf{Opérations réalisées} & \textbf{Données recueillies} \\
        \hline
        Irrigation et suivi hydrique & Arrosage régulier, observation du comportement du sol et adaptation des irrigations & Apports d'eau, temps d'infiltration, pluies et comportement du sol \\
        \hline
        Suivi phénologique & Observation de la germination, de la levée et de l'évolution des feuilles & Nombre de poquets levés, apparition des feuilles, déformations et changements de couleur \\
        \hline
        Entretien des cultures & Observation de l'état général des plants & Anomalies et mortalités observées \\
        \hline
        Suivi phytosanitaire & Observation des ravageurs, des dégâts et des symptômes de maladies ; application des traitements nécessaires & Ravageurs observés, symptômes, dégâts et traitements appliqués \\
        \hline
        Documentation & Enregistrement des observations dans l'agenda de culture et prise de photographies & Historique de l'évolution des cultures et des interventions \\
        \hline
    \end{tabular}
\end{table}

\subsection{Entretien des cultures}
L’entretien des cultures a été effectué régulièrement sur l’ensemble des six plots en fonction de leur état et du stade de développement des cultures. Il comprenait principalement \textbf{le désherbage, le binage, le sarclage, le sarclo-binage, la délimitation de cuvettes, le buttage, le tuteurage, l’irrigation, la fertilisation et la récolte} afin de maintenir des conditions de développement favorables aux cultures.

Le \textbf{désherbage et le sarclage} ont consisté à éliminer les adventices afin de limiter leur concurrence pour la lumière, l’eau, l’espace et les nutriments. Le \textbf{binage}, souvent associé au sarclage pour former le sarclo-binage, a permis d’ameublir la couche superficielle du sol, d’améliorer son aération et de faciliter l’infiltration de l’eau tout en réduisant son évaporation. La \textbf{délimitation des cuvettes} a favorisé la rétention de l’eau dans la zone racinaire et limité le ruissellement. Le \textbf{tuteurage} du gombo a servi de soutien aux plants pour réduire les risques de verse et de casse, tandis que le \textbf{buttage} du navet a contribué à protéger son tubercule et à favoriser son développement.
Certaines opérations d’entretien ont été adaptées aux stades de développement et aux exigences spécifiques de chaque culture. L’habillage du GOM TS a permis de réduire l’excès de feuillage et de favoriser la circulation de la lumière et de l’air. Le buttage du navet a été effectué lors de la formation du tubercule. Enfin, la récolte progressive a été réalisée selon le stade de maturité propre au gombo, au navet et à la laitue, permettant de suivre leur évolution et leur rendement.
\begin{figure}[H]
    \centering
    \imgstd{images/29_binage_sarclo_tuteurage.jpg}
    \caption{Actions respectives de Sarclo-binage, de Binage et de tuteurage. }
    \label{fig:entretient3}
\end{figure}

\subsection{Apport journalier d’eau}
Le pilotage de l’irrigation repose sur le suivi de l’humidité du sol et sur l’estimation des besoins en eau des cultures à partir de l’ETc calculée dans le classeur Excel.

\begin{table}[H]
    \centering
    \small
    \caption{Protocole de suivi et de pilotage de l'irrigation}
    \label{tab:pilotage_irrigation}
    \begin{tabular}{|p{4.2cm}|p{9.8cm}|}
        \hline
        \textbf{Éléments suivis} & \textbf{Modalités appliquées} \\
        \hline
        Plots concernés & GOM TS, NAV TS et LAI TS \\
        \hline
        Points de mesure & Deux points par plot \\
        \hline
        Outil de mesure & Capteur 8-en-1 \\
        \hline
        Fréquence de mesure & Deux fois par jour \\
        \hline
        Critère de déclenchement & Irrigation lorsque $S_a < S_c$ \\
        \hline
        Absence d'irrigation & Lorsque $S_a \geq S_c$ \\
        \hline
        Calcul de la dose & $ET_c$ calculée quotidiennement dans le classeur Excel à partir des données météorologiques collectées ou disponibles en prévision \\
        \hline
        Répartition de la dose & Deux apports : matin et soir \\
        \hline
        Durée d'application & Jusqu'à la fin de l'expérimentation \\
        \hline
    \end{tabular}
\end{table}

\begin{figure}[H]
    \centering
    \imgstd{images/30_collecte_hebdomadaire.jpg}
    \caption{Collecte des données sur le plot LAI TS. }
    \label{fig:collecte3}
\end{figure}

\subsection{Stratégie de lutte contre les maladies et ravageurs}
Au cours de notre expérimentation, divers ravageurs et maladies ont été observés sur le gombo, le navet et la laitue. Pour chaque cas, les symptômes ou dégâts, les agents responsables et les méthodes de lutte appliquées ont été identifiés et consignés.

\subsubsection{Ravageurs observés et méthodes de lutte}
Au cours de l’expérimentation, plusieurs ravageurs ont été observés sur les cultures, notamment le gombo et le navet. Les observations réalisées et les méthodes de lutte appliquées sont présentées dans le tableau suivant.

{\small
\begin{longtable}{|p{2.6cm}|p{3.4cm}|p{4.6cm}|p{4cm}|}
    \caption{Ravageurs observés, dégâts et méthodes de lutte}
    \label{tab:ravageurs_lutte} \\
    \hline
    \textbf{Ravageurs observés} & \textbf{Symptômes et dégâts constatés} & \textbf{Méthodes de lutte appliquées} & \textbf{Résultats observés} \\
    \hline
    \endfirsthead
    \hline
    \textbf{Ravageurs observés} & \textbf{Symptômes et dégâts constatés} & \textbf{Méthodes de lutte appliquées} & \textbf{Résultats observés} \\
    \hline
    \endhead
    Cochenilles farineuses (gombo) & Amas blancs et cotonneux sous les feuilles, jaunissement foliaire et présence de miellat. & Retrait manuel des cochenilles et utilisation ponctuelle d'un jet d'eau. & Les cochenilles sont éliminées des feuilles après l'intervention. Cependant, leur présence est de nouveau constatée quelques jours plus tard. \\
    \hline
    Cicadelle écumeuse (gombo) & Présence de mousse blanche sur les feuilles, protégeant les larves. & Jet d'eau pour éliminer la mousse et retrait manuel à l'éponge des larves et de la mousse. & Les larves sont délogées et éliminées lors du traitement. Toutefois, de nouvelles présences sont observées quelques jours après l'intervention. \\
    \hline
    Chenilles défoliatrices (gombo et navet) & Trous dans les feuilles et défoliation. & Retrait manuel et pulvérisation de solutions à base de neem ou d'ail, généralement en fin de journée. La solution de neem était préparée avec 500~g de feuilles pilées dans 10~L d'eau et une petite quantité de savon. La solution d'ail était préparée avec 500~g de gousses dans 9~L d'eau, complétée par une solution savonneuse. & L'odeur des solutions semble éloigner temporairement les chenilles, mais celles-ci réapparaissent après quelques jours. Lors des applications répétées d'une même solution, son effet répulsif semble diminuer progressivement. \\
    \hline
    Criquets (gombo et navet) & Consommation et sectionnement des tissus foliaires. & Pulvérisation d'une solution à base d'ail sur le feuillage. & La solution semble efficace au début, en éloignant les criquets. Toutefois, son effet répulsif diminue avec le temps, et les criquets finissent par revenir. \\
    \hline
\end{longtable}
}

Les ravageurs ont été identifiés à partir des observations effectuées sur les cultures. Les méthodes de lutte ont été adaptées aux types de dégâts constatés. Les photographies suivantes illustrent les ravageurs observés durant l’expérimentation.

    \begin{figure}[H]
    \centering
    \imgstd{images/31_cochenille_farineuse.jpg}
    \caption{Présence des cochenilles farineuses sur une feuilles de gombo. }
    \label{fig:cochenille3}
\end{figure}

    \begin{figure}[H]
    \centering
    \imgstd{images/32_cicadelle_ecumeuse.jpg}
    \caption{Cicadelle écumeuse sur une feuille de gombo : (a) mousse protectrice ; (b) larve extraite de la mousse. }
    \label{fig:cicadelle}
\end{figure}

\begin{figure}[H]
    \centering
    \imgstd{images/33_chenille.jpg}
    \caption{Ravageurs observés : 2 types de chenilles (en haut) et 2 types de criquets (en bas). }
    \label{fig:chenille3}
\end{figure}

\subsubsection{Maladies observées et méthodes de lutte}
L'observation a montré des symptômes compatibles avec ceux des maladies fongiques.

\begin{itemize}[label=$\bullet$]
   	\item \textbf{Les maladies fongiques : } sont causées par le développement de champignons pathogènes, accéléré par la combinaison d'une intense chaleur et d'une forte humidité. Trois types de maladies fongiques ont particulièrement été observés au cours de notre expérimentation : \textbf{la fusariose, le phytophthora et l'anthracnose.}

    \begin{itemize}[label=$\circ$]
   	\item \textbf{La fusariose}
ou flétrissement vasculaire, est principalement causée par un champignon appelé \textbf{Fusarium oxysporum}. Il pénètre l'organisme par les racines et colonise les vaisseaux de conduction de la sève (le xylème), sécrétant des toxines qui obstruent progressivement les canaux. Elle se manifeste par une chlorose et un flétrissement soudain, qui commencent généralement d'un côté de la plante ; une section de la tige révèle des vaisseaux internes brunis ou noircis. La sève ne circulant plus, la plante se dessèche en interne et meurt rapidement. Cette maladie est favorisée par un sol réchauffé (25 à 32°C), un pH acide et est encouragée par un stress hydrique causé par des périodes sèches alternées avec des inondations. De plus, le champignon a la capacité de survivre plusieurs années dans le sol, sous forme de spores résistantes.

    \item Le \textbf{phytophthora} responsable de la pourriture du collet et des racines, n'est pas un champignon à proprement parler mais un oomycète, un pathogène aquatique qui exploite un sol gorgé d'eau pour infiltrer les racines avec ses spores en mouvement. Il détruit les cellules racinaires afin d'absorber les matières organiques. Les signes de la maladie se manifestent par un brunissement et un ramollissement du collet, des racines noircies et visqueuses, et un flétrissement soudain des feuilles supérieures, entraînant l'effondrement et la mort de la plante en quelques jours. Cette maladie est favorisée par un excès d'eau (sols compacts, argileux ou trop irrigués), des températures tempérées à élevées  (15 à 25°C) et le manque d'oxygène au niveau des racines.

    \item \textbf{L'anthracnose} est une maladie foliaire et fruitière causée principalement par des champignons du genre \textbf{Colletotrichum}. Elle se maintient sur les résidus végétaux ou les graines et se répand par dispersion des spores lors des pluies. Elle se manifeste par des taches circulaires concaves, allant du brun au noir, parfois entourées d'un halo jaune sur les feuilles, ainsi que par des lésions noires et sèches sur les fruits (les capsules de gombo), qui les rendent non comestibles et non vendables. Une attaque sévère peut provoquer un dessèchement total et une chute des feuilles, engendrant une baisse significative du rendement. Elle est favorisée par une forte humidité (une pluie fréquente, une rosée matinale étendue, ou une irrigation par aspersion), une chaleur oppressante (20 à 28°C), ainsi qu'une densité de plantation excessive limitant la circulation de l'air.
    
    \item Méthode de lutte : Les plantes déjà atteintes n’ont pas pu être sauvées. Pour les plantes encore saines, une lutte préventive a été réalisée à l'aide de la solution d’ail et d’un traitement au bicarbonate, préparés à raison d’une cuillère à café par litre d’eau. Les traitements sont appliqués sur les feuilles en fin de journée (19h--20h).
	\end{itemize}
\end{itemize}

    \begin{figure}[H]
    \centering
    \imgstd{images/34_fongique.jpg}
    \caption{Gombo atteint d'une maladie fongique montrant un collet et des racines pourries. }
    \label{fig:fongique3}
\end{figure}

\subsection{Évolutions des différents plots}
Tout au long du cycle expérimental, les cultures ont fait l’objet d’une surveillance visuelle et qualitative quotidienne sur l'ensemble des plots, portant notamment sur \textbf{la vigueur, le feuillage et la croissance} des cultures. 
Ces observations ont été complétées par des photographies, des mesures régulières et le suivi des stades phénologiques clés, permettant une comparaison entre les plots T et TS. Les attaques de ravageurs et les maladies fongiques ont également été surveillées et consignées dans notre carnet de suivi expérimental.



















\section{Conclusion} 
Ce chapitre a présenté la démarche méthodologique mise en œuvre depuis l’installation des plots jusqu’au suivi des cultures, en passant par le pilotage de l’irrigation, la collecte des données, l’entretien et la lutte contre les ravageurs et les maladies. Cette approche a permis de mettre en place une base de données issue directement du terrain, regroupant ainsi les paramètres climatiques, hydriques et physico-chimiques nécessaires au calcul de l'évapotranspiration et au pilotage de l’irrigation. Ces données réelles constituent le socle du développement de l’application, destiné à automatiser le  calcul de l'évapotranspiration, et le pilotage de l’irrigation, tout en restant adaptée aux conditions de notre milieu d'étude.

\chapter{ MODÉLISATION, CONCEPTION ET IMPLÉMENTATION DU SYSTÈME}
\section{Introduction}
Le chapitre précédent a présenté la mise en place des plots et des cultures, le pilotage de l'irrigation selon la méthode FAO-56 Penman-Monteith \cite{allen_crop_1998-1}, les outils d'acquisition ainsi que la gestion hebdomadaire et le suivi des cultures. 

Cette expérimentation a permis d'établir une base de données issue du terrain, destinée à servir de fondement à la conception d’une application capable d’automatiser ces calculs et d’appuyer le pilotage de l’irrigation. 
Ce chapitre marque donc une transition de l'expérimentation agronomique vers une conception informatique du système.

La conception du système nécessite d’abord de déterminer les éléments à automatiser, les acteurs concernés et la logique de fonctionnement. 
Pour structurer cette analyse, la méthode Merise a été retenue. Ainsi, ce chapitre présente d’abord les raisons de ce choix, puis le domaine d'étude et ses acteurs, ensuite les différents modèles de conception, notamment le MCC, le MCT et le MOT ainsi que le diagramme de Gantt. Enfin, il présente l’application issue de ces travaux.

\section{Choix de la méthode d’analyse : Pourquoi MERISE}
Plusieurs méthodes d'analyse et de conception des systèmes d'information existent : Merise (Méthode d’Étude et de Réalisation Informatique pour les Systèmes d’Entreprise), UML (Unified Modeling Language), ou encore des approches agiles qui privilégient des itérations rapides plutôt que toute modélisation conceptuelle en amont. Le choix d'une méthode n'est jamais neutre : il détermine comment le problème est découpé, ce qui est formalisé en premier, et ce qui passe au second. Quatre raisons principales nous ont amenés à choisir Merise, chacune directement liée à la nature particulière de notre travail.

\subsection{Une méthode qui sépare la réflexion de la technique}
La méthode Merise  distingue trois niveaux d’abstraction : le niveau conceptuel(QUOI ?) qui définit les acteurs, les informations et les traitements ; le niveau Organisationnel (QUI? QUAND ? OÙ ?), qui précise la répartition des tâches ; et le niveau physique ou opérationnel (COMMENT ?), qui concerne les choix techniques. Cette démarche correspond à notre approche : définir et valider le fonctionnement du système à partir de l'expérimentation avant de penser aux solutions informatiques. Ainsi, Merise permet de formaliser le processus étudié avant sa mise en œuvre technique.

\subsection{Une méthode adaptée à un processus déjà éprouvé sur le terrain}
Contrairement à une démarche de conception qui part principalement des besoins prévisionnels des utilisateurs, notre travail s’appuie sur un processus agricole effectivement mis en œuvre et observé au cours de l’expérimentation. Les rôles des acteurs, les informations échangées et l’enchaînement des opérations ont ainsi pu être identifiés à partir des pratiques réalisées sur le terrain. Il s’agit notamment de l’étude du sol, du diagnostic, de la correction si nécessaire, du suivi de l'évolution des plots, du calcul des besoins en eau, de la prise de décision, de l’intervention et du compte rendu.

Le choix de Merise permet de structurer cette analyse en distinguant les échanges d’informations, les traitements réalisés et leur organisation opérationnelle, notamment à travers le MCC, le MCT et le MOT. UML aurait également pu être mobilisée pour représenter certains aspects du processus, mais Merise a été retenue en raison de sa cohérence avec la démarche de modélisation adoptée dans ce mémoire.

\subsection{Une méthode qui n’impose pas de choix technique prématuré et pédagogiquement adaptée et largement éprouvée}
Merise permet de décrire les traitements et de préciser leurs modalités d’exécution dans le MOT. Cette distinction est pertinente dans notre étude, où certaines opérations, comme la connexion du capteur 8-en-1, restent manuelles, tandis que d’autres, notamment les calculs, sont destinées à être automatisées dans l’application envisagée.

Le choix de Merise s’inscrit également dans le cadre de notre formation à l’USSEIN et correspond aux besoins de modélisation de cette étude. UML aurait également pu être utilisée, mais Merise a été retenue pour sa cohérence avec la démarche adoptée.

\subsection{Synthèse comparative}
Ce tableau récapitule les critères qui ont orienté ce choix méthodologique.

{\small
\begin{longtable}{|p{3.2cm}|p{5.4cm}|p{5.4cm}|}
    \caption{Critères orientant le choix méthodologique}
    \label{tab:criteres_choix_methodologique} \\
    \hline
    \textbf{Critères} & \textbf{Merise} & \textbf{UML} \\
    \hline
    \endfirsthead
    \hline
    \textbf{Critères} & \textbf{Merise} & \textbf{UML} \\
    \hline
    \endhead
    Point de départ de la modélisation & Permet de structurer les flux d'informations, les traitements et leur organisation à partir du processus observé. & Permet de représenter les besoins, les interactions entre acteurs et le comportement du système, notamment à travers les cas d'utilisation et les diagrammes de comportement. \\
    \hline
    Séparation entre les aspects conceptuels et organisationnels ou techniques & Distingue plusieurs niveaux de modélisation, notamment à travers le MCT et le MOT. & Propose différents diagrammes pour représenter les aspects fonctionnels, structurels et comportementaux du système. \\
    \hline
    Représentation des opérations manuelles et automatisées & Permet de préciser les modalités d'exécution des traitements dans le modèle organisationnel. & Permet également de représenter les interactions humaines et les traitements automatisés à l'aide de diagrammes appropriés. \\
    \hline
    Adéquation au processus étudié & Adaptée à la formalisation du processus agricole observé et à la distinction entre traitements et organisation. & Également envisageable, notamment pour représenter les interactions entre les utilisateurs et l'application future. \\
    \hline
    Cohérence avec la formation et le cadre du mémoire & Correspond à la démarche de modélisation retenue et au cadre méthodologique de la formation. & Constitue une autre approche possible, qui aurait nécessité une organisation différente de la modélisation. \\
    \hline
\end{longtable}
}

En définitive, Merise a été retenue pour sa cohérence avec le processus étudié et les objectifs de cette étude, sans exclure la pertinence d’UML pour d’autres approches de modélisation.



\section{Du terrain au système d’information : Domaine d'étude et acteur}
Avant la présentation des modèles conceptuels, il est nécessaire de définir le domaine d'étude et d'identifier les acteurs. En Merise, un acteur est une personne, entité, organisation, dispositif ou système externe échangeant des informations avec le domaine, il peut être interne ou externe. Dans notre cas, ces acteurs ont été identifiés à partir des rôles réellement observés et exercés lors de l'expérimentation, permettant ainsi d'établir leur correspondance avec le système d’information.

\begin{longtable}{|p{7cm}|p{3.8cm}|p{2.2cm}|}
    \caption{Acteurs du système d'information, avec leur type et leur rôle sur le terrain}
    \label{tab:acteurs_si} \\
    \hline
    \textbf{Rôle observé sur le terrain (chapitre 3 et protocole)} & \textbf{Acteur du système d'information} & \textbf{Type} \\
    \hline
    \endfirsthead
    \hline
    \textbf{Rôle observé sur le terrain (chapitre 3 et protocole)} & \textbf{Acteur du système d'information} & \textbf{Type} \\
    \hline
    \endhead
    Équipe expérimentale, relevant les données du sol à l'aide du capteur 8-en-1 deux fois par jour (matin et soir), calculant $ET_0$/$ET_c$, décidant d'irriguer ou non, appliquant les dosages d'amendement (protocole) & Technicien & Interne \\
    \hline
    Équipe expérimentale, finançant l'achat du capteur 8-en-1 et des composants de l'amendement (biochar, coque d'arachide, compost, biofertilisant) & Bailleur & Interne \\
    \hline
    Responsable scientifique & Utilisateur & Externe \\
    \hline
    Le capteur 8-en-1, mesurant humidité, température, pH, conductivité électrique, N, P, K & Capteur 8-en-1 & Interne \\
    \hline
    L'application météo du smartphone (AccuWeather) utilisée pendant l'expérimentation, remplacée dans l'application par l'API OpenWeather & AccuWeather & Externe \\
    \hline
    Le fournisseur du capteur 8-en-1 & Fournisseur capteur & Externe \\
    \hline
    Le fournisseur des composants de l'amendement (charbon de bois, coque d'arachide, compost, engrais bio) & Fournisseur intrant & Externe \\
    \hline
\end{longtable}

Cette correspondance entre l'expérience vécue et les acteurs modélisés assure que le système conçu automatise une pratique réelle documentée et testée, basée sur les conditions agronomiques et climatiques observées sur le champ expérimental Agro TIC de l'USSEIN.

\section{Les deux grandes catégories de modèles et le diagramme de Gantt}
En Merise, on distingue deux grandes catégories de modèles : les modèles de traitement (MCC, MCT et MOT) et les modèles de données (MCD et MLD). Dans notre méthodologie de conception, nous avons étudié les trois différents modèles de traitements.
La représentation graphique des modèles est presentée a la fin de chaque sous section concernée et le diagramme de Gantt sera présenté dans l’Annexe~\ref{ann:diagrammes}.

\subsection{Modèle Conceptuel de Communication (MCC)}
Le MCC représente les acteurs impliqués dans le processus étudié ainsi que les flux d’informations échangés entre eux. 
Dans cette étude, il met en évidence les relations entre les acteurs internes et externes, notamment \textbf{le technicien, l’utilisateur, les fournisseurs, le bailleur et le capteur 8-en-1}. 

\begin{longtable}{|p{1cm}|p{6.5cm}|p{6.5cm}|}
    \caption{Représentation des flux}
    \label{tab:flux} \\
    \hline
    \textbf{N°} & \textbf{Émetteur $\rightarrow$ Récepteur} & \textbf{Libellé du flux} \\
    \hline
    \endfirsthead
    \hline
    \textbf{N°} & \textbf{Émetteur $\rightarrow$ Récepteur} & \textbf{Libellé du flux} \\
    \hline
    \endhead
    1 & Technicien $\rightarrow$ Fournisseur capteur & Commande de capteurs \\
    \hline
    2 & Fournisseur capteur $\rightarrow$ Technicien & Établissement et transmission du devis \\
    \hline
    3 & Technicien $\rightarrow$ Fournisseur capteur & Devis accepté \\
    \hline
    4 & Technicien $\rightarrow$ Bailleur & Demande de fonds (capteur) \\
    \hline
    5 & Bailleur $\rightarrow$ Technicien & Fonds mis à disposition (capteur) \\
    \hline
    6 & Technicien $\rightarrow$ Fournisseur capteur & Paiement \\
    \hline
    7 & Fournisseur capteur $\rightarrow$ Technicien & Livraison \\
    \hline
    8 & Technicien $\rightarrow$ Fournisseur intrant & Commande des intrants \\
    \hline
    9 & Fournisseur intrant $\rightarrow$ Technicien & Établissement et transmission du devis \\
    \hline
    10 & Technicien $\rightarrow$ Fournisseur intrant & Devis accepté \\
    \hline
    11 & Technicien $\rightarrow$ Bailleur & Demande de fonds (intrants) \\
    \hline
    12 & Bailleur $\rightarrow$ Technicien & Fonds mis à disposition (intrants) \\
    \hline
    13 & Technicien $\rightarrow$ Fournisseur intrant & Paiement \\
    \hline
    14 & Fournisseur intrant $\rightarrow$ Technicien & Livraison \\
    \hline
    15 & Technicien $\rightarrow$ Capteur 8-en-1 & Configuration du capteur \\
    \hline
    16 & Capteur 8-en-1 $\rightarrow$ Technicien & Envoi des données du sol \\
    \hline
    17 & Technicien $\rightarrow$ OpenWeather & Requête \\
    \hline
    18 & OpenWeather $\rightarrow$ Technicien & Envoi des données météorologiques \\
    \hline
    19 & Technicien $\rightarrow$ Utilisateur & Transmission des résultats et recommandations \\
    \hline
    20 & Utilisateur $\rightarrow$ Technicien & Retour \\
    \hline
\end{longtable}

\begin{figure}[H]
    \centering
    \includegraphics[width=\textwidth]{images/MCC.pdf}
    \caption{Modèle conceptuel de communication (MCC). }
    \label{fig:mcc}
\end{figure}

\subsection{Modèle Conceptuel de Traitement (MCT)} 
Le MCT décrit l’enchaînement des opérations du processus, depuis l’identification des besoins en équipements jusqu’à la transmission des résultats. 
Il met en évidence les \textbf{événements déclencheurs, les opérations réalisées, les conditions de leur exécution et les résultats obtenus}. 
Le tableau récapitulatif présente les principales opérations du processus.

{\small
\begin{longtable}{|p{0.8cm}|p{4.2cm}|p{4.2cm}|p{4.8cm}|}
    \caption{Correspondance entre les opérations, les événements déclencheurs et les résultats du MCT}
    \label{tab:operations_mct} \\
    \hline
    \textbf{N°} & \textbf{Opération} & \textbf{Événement(s) déclencheur(s)} & \textbf{Résultat} \\
    \hline
    \endfirsthead
    \hline
    \textbf{N°} & \textbf{Opération} & \textbf{Événement(s) déclencheur(s)} & \textbf{Résultat} \\
    \hline
    \endhead
    1 & Commander le capteur 8-en-1 auprès du fournisseur capteur & Besoin du capteur identifié et accord de l'utilisateur & En attente du devis \\
    \hline
    2 & Traitement du devis & Devis reçu & Acceptation ou refus ; si refus, fin de prestation \\
    \hline
    3 & Paiement & Acceptation du devis, demande de fonds et réception des fonds & Paiement effectué \\
    \hline
    4 & Installer et configurer le capteur & Capteur livré (ou reçu) & Capteur opérationnel \\
    \hline
    5 & Réaliser l'étude du sol (collecte des données du sol à l'aide du 8-en-1) & Accord de l'utilisateur & Collecte de données effectuée \\
    \hline
    6 & Analyse des résultats & Collecte effectuée & Besoin de correction identifié (si résultats valides) et ne pas se prononcer (si résultats invalides) \\
    \hline
    7 & Commande d'intrants & Accord de l'utilisateur & Attente du devis \\
    \hline
    8 & Traitement du devis & Devis reçu & Fin de prestation si devis refusé, ou demande de fonds si devis accepté \\
    \hline
    9 & Paiement & Fonds reçus & Paiement effectué \\
    \hline
    10 & Réceptionner les intrants & Livraison d'intrants & Intrants reçus \\
    \hline
    11 & Réaliser l'amendement sur la parcelle & Besoin de correction identifié et intrants reçus & Parcelle amendée \\
    \hline
    12 & Configuration et relevés de données (à chaque connexion du capteur) & Parcelle amendée & Données du sol reçues \\
    \hline
    13 & Requête AccuWeather & Parcelle amendée & Données météo reçues si connexion disponible, ou en attente de connexion si connexion non disponible \\
    \hline
    14 & Calculer les indicateurs agronomiques & Données du sol reçues et données météo reçues & Indicateurs agronomiques calculés \\
    \hline
    15 & Évaluer le besoin d'irrigation & Indicateurs agronomiques calculés & Irrigation nécessaire si $S_a < S_c$, ou stock suffisant (ne pas irriguer) si $S_a \geq S_c$ \\
    \hline
    16 & Irriguer la parcelle & Irrigation nécessaire et eau disponible & Irrigation réalisée \\
    \hline
    17 & Évaluer le besoin de fertilisation & Indicateurs agronomiques calculés & Fertilisation nécessaire si déficit de N/P/K, ou stock suffisant si pas de déficit \\
    \hline
    18 & Fertiliser la parcelle & Fertilisation nécessaire et disponibilité des intrants & Fertilisation réalisée \\
    \hline
    19 & Consolider les résultats et rédiger le rapport d'avancement & Irrigation réalisée, stock suffisant (ne pas irriguer), fertilisation réalisée et stock suffisant (ne pas fertiliser) & Rapport d'avancement rédigé \\
    \hline
    20 & Transmission du rapport & Rapport d'avancement rédigé & \\
    \hline
\end{longtable}
}

\begin{figure}[H]
    \centering
    \includegraphics[width=\textwidth, height=0.85\textheight, keepaspectratio]{images/MCT1.pdf}
    \caption{Modèle conceptuel de traitement (MCT), partie 1/6. }
    \label{fig:mct1}
\end{figure}

\begin{figure}[H]
    \centering
    \includegraphics[width=\textwidth, height=0.85\textheight, keepaspectratio]{images/MCT2.pdf}
    \caption{Modèle conceptuel de traitement (MCT), partie 2/6. }
    \label{fig:mct2}
\end{figure}

\begin{figure}[H]
    \centering
    \includegraphics[width=\textwidth, height=0.85\textheight, keepaspectratio]{images/MCT3.pdf}
    \caption{Modèle conceptuel de traitement (MCT), partie 3/6. }
    \label{fig:mct3}
\end{figure}

\begin{figure}[H]
    \centering
    \includegraphics[width=\textwidth, height=0.85\textheight, keepaspectratio]{images/MCT4.pdf}
    \caption{Modèle conceptuel de traitement (MCT), partie 4/6. }
    \label{fig:mct4}
\end{figure}

\begin{figure}[H]
    \centering
    \includegraphics[width=\textwidth, height=0.85\textheight, keepaspectratio]{images/MCT5.pdf}
    \caption{Modèle conceptuel de traitement (MCT), partie 5/6. }
    \label{fig:mct5}
\end{figure}

\begin{figure}[H]
    \centering
    \includegraphics[width=\textwidth, height=0.85\textheight, keepaspectratio]{images/MCT6.pdf}
    \caption{Modèle conceptuel de traitement (MCT), partie 6/6. }
    \label{fig:mct6}
\end{figure}

\subsection{Modèle Organisationnel de Traitement (MOT)}
Le MOT complète le MCT en précisant, pour chaque opération, les modalités organisationnelles de son exécution, notamment \textbf{le lieu, le moment et la nature} du traitement. 
Il permet ainsi de distinguer les opérations manuelles de celles envisagées comme automatisables.

\begin{figure}[H]
    \centering
    \includegraphics[width=\textwidth, height=0.85\textheight, keepaspectratio]{images/MOT1_1.pdf}
    \caption{Modèle organisationnel de traitement (MOT), partie 1/9. }
    \label{fig:mot1.1}
\end{figure}

\begin{figure}[H]
    \centering
    \includegraphics[width=\textwidth, height=0.85\textheight, keepaspectratio]{images/MOT1_2.pdf}
    \caption{Modèle organisationnel de traitement (MOT), partie 2/9. }
    \label{fig:mot1.2}
\end{figure}

\begin{figure}[H]
    \centering
    \includegraphics[width=\textwidth, height=0.85\textheight, keepaspectratio]{images/MOT1_3.pdf}
    \caption{Modèle organisationnel de traitement (MOT), partie 3/9. }
    \label{fig:mot1.3}
\end{figure}

\begin{figure}[H]
    \centering
    \includegraphics[width=\textwidth, height=0.85\textheight, keepaspectratio]{images/MOT1_4.pdf}
    \caption{Modèle organisationnel de traitement (MOT), partie 4/9. }
    \label{fig:mot1.4}
\end{figure}

\begin{figure}[H]
    \centering
    \includegraphics[width=\textwidth, height=0.85\textheight, keepaspectratio]{images/MOT1_5.pdf}
    \caption{Modèle organisationnel de traitement (MOT), partie 5/9. }
    \label{fig:mot1.5}
\end{figure}

\begin{figure}[H]
    \centering
    \includegraphics[width=\textwidth, height=0.85\textheight, keepaspectratio]{images/MOT1_6.pdf}
    \caption{Modèle organisationnel de traitement (MOT), partie 6/9. }
    \label{fig:mot1.6}
\end{figure}

\begin{figure}[H]
    \centering
    \includegraphics[width=\textwidth, height=0.85\textheight, keepaspectratio]{images/MOT1_7.pdf}
    \caption{Modèle organisationnel de traitement (MOT), partie 7/9. }
    \label{fig:mot1.7}
\end{figure}

\begin{figure}[H]
    \centering
    \includegraphics[width=\textwidth, height=0.85\textheight, keepaspectratio]{images/MOT1_8.pdf}
    \caption{Modèle organisationnel de traitement (MOT), partie 8/9. }
    \label{fig:mot1.8}
\end{figure}

\begin{figure}[H]
    \centering
    \includegraphics[width=\textwidth, height=0.85\textheight, keepaspectratio]{images/MOT1_9.pdf}
    \caption{Modèle organisationnel de traitement (MOT), partie 9/9. }
    \label{fig:mot1.9}
\end{figure}
\section{Conception de l’application AgriSens}
Les modèles présentés dans ce chapitre ont servi de base à la conception de l'application, une application web de type Progressive Web App (PWA) a été conçue afin d’automatiser le suivi et le pilotage de l'irrigation de précision et de la fertilisation, en accord avec le procédé approuvé au chapitre 3.

\subsection{Architecture technique}
L'application repose sur la plateforme FIWARE, adaptée à la gestion des données contextuelles issues des capteurs et des sources météorologiques identifiées dans le MCC. L’architecture retenue s’appuie sur les composants ci-dessous : 

\textit{Orion Context Broker (protocole NGSI-v2)} : le cœur du système, qui stocke et distribue les entités du domaine (parcelles, relevés du capteur) sous forme d'informations contextuelles standardisées.

\textit{MongoDB} est la base de données sous-jacente d'Orion.

\textit{Keycloak} est le gestionnaire d'identité et d'accès, chargé de l'authentification des utilisateurs et de l'attribution de leurs rôles applicatifs (Administrateur, Agronome, Technicien).

\textit{Wilma, proxy PEP (Policy Enforcement Point)}, s’intercale entre l'interface web et Orion pour vérifier, en interrogeant Keycloak, que chaque requête est authentifiée avant d'atteindre le Context Broker, et complété par Wilma-auth, qui gère l'inscription des comptes.

Cette architecture donne un contenu technique concret aux flux 15 et 16 du MCC (« Configuration du capteur » et « Envoi des données du sol ») : dans l'application, ces échanges sont représentés par des requêtes NGSI-v2 vers Orion, toujours sécurisées par Wilma.

\subsection{Organisation de l’application et rôles}
L'application distingue trois rôles applicatifs, qui affinent l'acteur unique « Technicien » du MCC en fonction du niveau d'accès réellement nécessaire, comme le montre le tableau ci-dessous.

\begin{table}[H]
    \centering
    \small
    \caption{Correspondance entre les rôles applicatifs, les accès et les acteurs conceptuels}
    \label{tab:roles_applicatifs}
    \begin{tabular}{|p{2.4cm}|p{5.2cm}|p{6.4cm}|}
        \hline
        \textbf{Rôle applicatif} & \textbf{Accès} & \textbf{Correspondance avec les acteurs conceptuels (MCC/MCT)} \\
        \hline
        Technicien & Accueil, Parcelles, Capteurs, Calculs, Graphes, Historique, plus gestion des utilisateurs et paramétrage des cultures ($K_c$, seuils N/P/K) & Acteur \og Technicien \fg{} : relève les données, déclenche les calculs, agit sur la parcelle. Rôle transverse, nécessaire pour maintenir les référentiels agronomiques du protocole \\
        \hline
        Utilisateur & Identique, avec lecture consolidée des résultats & Proche de l'acteur \og Utilisateur \fg{} : observe les résultats \\
        \hline
    \end{tabular}
\end{table}

L'application s'organise en six pages destinées aux rôles Technicien et Agronome : Accueil, Parcelles, Capteurs, Calculs, Graphes et Historique auxquelles s'ajoute une septième page, exclusivement réservée à l'administrateur. Ce tableau établit la correspondance entre chacune de ces pages et l'opération du MCT qu'elle automatise.

\begin{table}[H]
    \centering
    \small
    \caption{Correspondance entre les pages de l'application et les opérations automatisées du MCT}
    \label{tab:pages_application}
    \begin{tabular}{|p{3.4cm}|p{10.2cm}|}
        \hline
        \textbf{Page de l'application} & \textbf{Opération(s) du MCT automatisée(s)} \\
        \hline
        Accueil (tableau de bord) & Vue de synthèse des indicateurs calculés et de l'état des parcelles \\
        \hline
        Parcelles & Gestion des entités correspondant aux plots test et témoin \\
        \hline
        Capteurs & \og Relever les données du sol \fg{} : lecture du capteur 8-en-1 via le protocole Modbus RTU réel, avec une stabilisation de 30 secondes avant enregistrement \\
        \hline
        Calculs & \og Calculer les indicateurs agronomiques \fg{} ($ET_0$, $ET_c$, $RU$, $RFU$, $S_c$, $S_a$, déficits N/P/K) et \og Évaluer le besoin d'irrigation/de fertilisation \fg{}. Seuls les résultats sont affichés, sans exposer les formules intermédiaires \\
        \hline
        Graphes & Historisation des relevés du capteur et visualisation de leur évolution \\
        \hline
        Historique & Traçabilité des interventions d'irrigation et de fertilisation \\
        \hline
        Technicien / Admin & Paramétrage par culture des coefficients $K_c$ (initial, mi-saison, final) et des seuils N/P/K optimaux \\
        \hline
    \end{tabular}
\end{table}

Conformément au MCT, l’application conserve les actions demandant une  interventions humaines, telles que le branchement physique du capteur 8-en-1 et la confirmation de lecture. Par contre, la chaîne de traitement est ensuite automatisée à partir des données du capteur et de OpenWeather, pour réduire les erreurs liées aux traitements manuels. La conception repose ainsi sur une automatisation respectant les pratiques réellement observées sur le terrain.


\section{Implémentation de l’application AgriSens}
L’application est directement accessible depuis un navigateur (Edge ou Chrome) et peut également être téléchargée et installée. Cette section présente son interface ainsi qu’un aperçu synthétique de ses principales fonctionnalités.
%%%%%%%%%% CREATION DE COMPTE : INSCRIPTION %%%%%%%%%%%%%

Lors du premier accès à l’application, la création de compte est obligatoire. La figure ci-dessus présente les étapes à remplir. 
Cette étape s’effectue une seule fois et uniquement lorsque l’utilisateur ne possède pas de compte. Ainsi, après le remplissage du formulaire, l’utilisateur clique sur le bouton \textit{“Créer mon compte”} afin de créer son compte.

\begin{figure}[H]
    \centering
    \begin{minipage}{0.45\textwidth}
        \centering
        \includegraphics[width=\textwidth]{images/1_a_pour_appli.jpeg}
    \end{minipage}
    \hfill
    \begin{minipage}{0.45\textwidth}
        \centering
        \includegraphics[width=\textwidth]{images/1_b_pour_appli.jpeg}
    \end{minipage}
    \caption{Interface de création de compte (parties 1 et 2). }
    \label{fig:inscription}
\end{figure}

Une fois le compte créé, l'utilisateur pourra enfin se connecter. Les figures en paralélles montrent la page de connexion à gauche et la page d'acceuil à droite.

%%%%%%% SE CONNECTER AVEC LE COMPTE CREER %%%%%%%%%%

\begin{figure}[H]
    \centering
    \begin{minipage}{0.45\textwidth}
        \centering
        \includegraphics[width=\textwidth]{images/1_c_pour_appli.jpeg}
    \caption{Interface de connexion. }
    \label{fig:conexion}
    \end{minipage}
    \hfill
    \begin{minipage}{0.45\textwidth}
        \centering
        \includegraphics[width=\textwidth]{images/1_d_pour_appli.jpeg}
    \caption{Interface de connexion. }
    \label{fig:dashboard}
    \end{minipage}
\end{figure}

Sur cette interface d'accueil figurent les informations récentes de l’utilisateur s’il s’est déjà connecté, dans le cas contraire, elle est vide. En dessous de cette interface, on retrouve les autres fonctionnalités de l'application : \textit{PARCELLES, CAPTEURS, CALCULS, GRAPHES et HISTORIQUES}.

{\small
\begin{longtable}{|p{2.6cm}|p{5.7cm}|p{5.7cm}|}
    \caption{Fonctionnalités de l'application, entrées et sorties}
    \label{tab:fonctionnalites_appli} \\
    \hline
    \textbf{Fonctionnalité} & \textbf{Entrées} & \textbf{Sorties} \\
    \hline
    \endfirsthead
    \hline
    \textbf{Fonctionnalité} & \textbf{Entrées} & \textbf{Sorties} \\
    \hline
    \endhead
    \textbf{Parcelles} & Nom et superficie de la parcelle, culture choisie, date de semis & Parcelle enregistrée, stade de développement et nombre de jours calculés automatiquement \\
    \hline
    \textbf{Capteurs} & Données météorologiques collectées automatiquement et données du sol collectées à l'aide du capteur 8-en-1 & Affichage des données météorologiques et des paramètres du sol \\
    \hline
    \textbf{Calculs} & Données météorologiques, données du sol et paramètres culturaux & $ET_0$, $ET_c$, dose d'irrigation ($D_i$), stock actuel d'eau ($S_a$) et décision d'irrigation \\
    \hline
    \textbf{Graphes} & Données de croissance des plantes, paramètres météorologiques et données d'irrigation & Courbes d'évolution des plantes, des paramètres météorologiques, des besoins hydriques et des doses d'irrigation apportées \\
    \hline
    \textbf{Historiques} & Données enregistrées au cours du cycle cultural & Historique de la spéculation choisie en fin de cycle \\
    \hline
\end{longtable}
}
\section{Conclusion}
Ce chapitre a permis de traduire le processus expérimental présenté au chapitre précédent en un système d’information structuré. Le choix de la méthode Merise a permis de représenter les acteurs, les flux d'informations et les enchaînements des traitements à travers le MCC, le MCT et le MOT. Cette modélisation a ainsi permis de conserver la logique du processus réellement observée sur le terrain, tout en identifiant les opérations pouvant être automatisées.

Sur cette base, l’application Agri Sens a été conçue sous la forme d’une progressive Web APP, en s’appuyant sur l’architecture FIWARE pour la gestion des données provenant du capteur 8-en-1 et les données provenant de la météo. L’application automatise notamment le calcul des indicateurs agronomiques, l'évaluation des besoins en fertilisation et en irrigation, ainsi que la visualisation des données et leur historique tout en conservant les interventions nécessitant une action humaine.

Ainsi, ce chapitre établit le passage entre le processus agronomique expérimenté et sa traduction informatique, en proposant une conception fidèle aux pratiques observées et aux besoins réels identifiés durant l'expérimentation.


\chapter{RÉSULTATS ET DISCUSSIONS}
\section{Introduction}
Les chapitres précédents ont permis de poser l'ensemble des bases de cette recherche : le chapitre 1 a présenté l'état de l'art, 
le chapitre 2 a présenté le protocole expérimental, le chapitre 3 a détaillé la mise en œuvre concrète de ce protocole sur le terrain, tandis que le chapitre 4 a présenté la modélisation et la conception de l'application de pilotage de l'irrigation et de la fertilisation.

Ce chapitre  rassemble et discute l'ensemble des résultats obtenus au cours de l'expérimentation, qu'ils soient issus du suivi du sol, 
de la croissance des cultures, du comportement phytosanitaire des plots, du rendement final, ou du fonctionnement de l'application en question. 
Il permet une interprétation honnête et nuancée de l'effet réel de l'amendement du sol et du pilotage de l'irrigation et de la fertilisation sur les performances agronomiques des trois cultures étudiées : le gombo, le navet et la laitue.

Le chapitre se répartit  en cinq grandes sections. En premier lieu, les résultats relatifs à l'amendement du sol sont présentés (section 5.2) ; ensuite, viennent les résultats des courbes d'évolution des paramètres étudiés (section 5.3) ; 
puis, en section 5.4, nous avons les résultats phytosanitaires ; enfin, les résultats des rendements sont aussi présentés en section 5.5 avant une conclusion partielle (section 5.6) menant aux recommandations et perspectives de ce travail.

\section{Résultats relatifs à l'amendement du sol}
Comme le décrivent nos chapitres précédents, pour chaque amendement utilisé, nous avons des attentes spécifiques. Cette section décrit d’abord les résultats attendus selon la littérature scientifique avant de présenter les résultats réellement obtenus.

\subsection{Mécanisme d’action attendu des composants de l’amendement}
Les différents composants de l’amendement ont été associés en raison de leurs effets complémentaires attendus sur le sol et les cultures. \textbf{Le biochar} contribue à la rétention des nutriments et à l’amélioration des propriétés du sol, tandis que le \textbf{compost} favorise sa fertilité et sa capacité de rétention en eau. La \textbf{PCA} constitue une source de matière organique à décomposition lente et contribue aussi à l'aération du sol. Enfin, \textbf{l'EGB}, grâce aux micro-organismes bénéfiques qu’il contient, est susceptible de favoriser la nutrition des plantes et leur tolérance au stress salin.

Ces effets constituent les mécanismes attendus dont la manifestation doit être appréciée à partir des observations réalisées au cours de l’expérimentation.

\subsection{Résultats des tests au début de l'expérimentation}
La présente section est consacrée à l’interprétation des données obtenues selon trois axes complémentaires : 
l’évaluation du comportement du capteur, la caractérisation des composantes de l’amendement utilisé et l’analyse comparative de l’état du sol entre les plots T et les plots TS.

\subsubsection{Validation du comportement du capteur}
Avant d’utiliser les données du sol collectées sur le terrain, la cohérence de la réponse du capteur a été vérifiée à différents niveaux d’humidité. Un même échantillon de sol non amendé a été réparti dans trois récipients correspondant à trois niveaux d'humidité différents : \textbf{sol sec, moitié humide et sol mouillé}. Les mesures obtenues de ces pots ont été relevées et leurs valeurs sont présentées dans le tableau ci-dessous :

\begin{table}[H]
    \centering
    \caption{Résultats sur les échantillons des pots}
    \label{tab:collecte_pots}
    \begin{tabular}{|p{2.4cm}|p{1.6cm}|p{1.9cm}|p{0.8cm}|p{1.9cm}|p{1.7cm}|p{1.7cm}|p{1.7cm}|}
        \hline
        \textbf{État du sol} & \textbf{Humidité (\%)} & \textbf{Température (°C)} & \textbf{pH} & \textbf{EC (\textmu S/cm)} & \textbf{N (mg/kg)} & \textbf{P (mg/kg)} & \textbf{K (mg/kg)} \\
        \hline
        Sol sec & 0--5 & 32,7--32,9 & 7 & 0 & 0 & 0 & 0 \\
        \hline
        Sol moitié humide & 7,2--7,8 & 31,9 & 7 & 126--193 & 6--9 & 8--13 & 20--30 \\
        \hline
        Sol mouillé & 100 & 31,3 & 10 & 5451--5570 & 272--278 & 381--389 & 872--891 \\
        \hline
    \end{tabular}
\end{table}

Les résultats montrent une augmentation régulière de $EC$, du $N$, du $P$ et du $K$ avec l'humidité du sol. Cette tendance s’explique par le fait que l’eau favorise la dissolution et le déplacement des ions, facilitant ainsi leur détection par le capteur. À l’inverse, un sol sec limite fortement leur déplacement, ce qui explique les valeurs nulles observées dans le pot de sol sec. La valeur de 100\% d'humidité mesurée dans le sol mouillé s’explique par sa saturation uniforme contrairement aux plots au champ ou le drainage naturel limite une telle saturation. Ce test en pot constitue donc principalement une vérification de la cohérence de la réponse du capteur à des niveaux d'humidité différents.

\begin{figure}[H]
    \centering
    \imgstd{images/27_echantillons_de_sol.jpg}
    \caption{Échantillons de sol à différents niveaux d'humidité. }
    \label{fig:echantillons3}
\end{figure}

\subsubsection{Caractérisation des composantes de l’amendement}
L'amendement anti-sel est constitué de compost, de biochar, de PCA et de EGB. Afin d’identifier la contribution de chacun, un test d’hydratation individuel a été effectué avant leur incorporation au sol. Les résultats obtenus sont présentés dans le tableau ci-dessous :

\begin{table}[H]
    \centering
    \small
    \caption{Résultats sur les composants de l'amendement}
    \label{tab:releves_amendement}
    \begin{tabular}{|p{2.4cm}|p{1.6cm}|p{1.9cm}|p{0.9cm}|p{1.9cm}|p{1.5cm}|p{1.6cm}|p{1.6cm}|}
        \hline
        \textbf{Composants} & \textbf{Humidité (\%)} & \textbf{Température (°C)} & \textbf{pH} & \textbf{EC (\textmu S/cm)} & \textbf{N (mg/kg)} & \textbf{P (mg/kg)} & \textbf{K (mg/kg)} \\
        \hline
        Biofertilisant & 42,9--43,6 & 31 & 6,9 & 1394--1397 & 69 & 97 & 223 \\
        \hline
        Biochar & 75,3--75,6 & 30,7--30,8 & 6,8--6,9 & 1605--1659 & 80--82 & 112--116 & 256--265 \\
        \hline
        Compost & 100 & 30,4 & 6,9--7,6 & 6112--6208 & 305--310 & 427--434 & 977--993 \\
        \hline
        Coque d'arachide & 52,2--52,3 & 30--31,5 & 6,9 & 1080--1085 & 54 & 75 & 172--173 \\
        \hline
    \end{tabular}
\end{table}

Ces résultats révèlent que le compost présente les valeurs les plus élevées de $EC$, de $N$, de $P$ et de $K$, en raison de sa richesse en MO minéralisée. Malgré cette EC forte, son rôle dans la lutte contre la salinité repose principalement sur l'amélioration de la structure, de la rétention et du drainage du sol favorisant ainsi le lessivage des sels excédentaires. Le EGB et la PCA caractérisés par les $EC$ les plus faibles participent quant à eux à la régulation de la salinité dans le mélange. La figure ci-dessous montre les échantillons analysés.

\begin{figure}[H]
    \centering
    \imgstd{images/28_prelevement_intrant.jpg}
    \caption{Prélèvements des 4 éléments de l’amendement. }
    \label{fig:prelevement3}
\end{figure}

\subsubsection{Analyse comparative de l’état du sol : plots témoins et plots test}
Après mélange et incorporation des quatres composants dans les plots TS suivi d’un arrosage identique, un relevé a été réalisé la veille du semis (le 24 juin 2026) afin de comparer l'état initial du sol entre plot TS amendé et plot T non amendé.

\begin{table}[H]
    \centering
    \small
    \caption{Résultats de l'analyse comparative entre les plots test et témoins}
    \label{tab:analyse_comparative}
    \begin{tabular}{|p{1.9cm}|p{1.6cm}|p{1.9cm}|p{0.9cm}|p{1.9cm}|p{1.5cm}|p{1.6cm}|p{1.6cm}|}
        \hline
        \textbf{Plot} & \textbf{Humidité (\%)} & \textbf{Température (°C)} & \textbf{pH} & \textbf{EC (\textmu S/cm)} & \textbf{N (mg/kg)} & \textbf{P (mg/kg)} & \textbf{K (mg/kg)} \\
        \hline
        GOM T & 79,3 & 32,3 & 10 & 2282,2 & 113,7 & 159,3 & 364,7 \\
        \hline
        GOM TS & 100 & 30,3 & 8,2 & 4615,6 & 230,3 & 322,6 & 738 \\
        \hline
        NAV T & 71,5 & 31,7 & 10 & 3010,8 & 150 & 210,4 & 481,4 \\
        \hline
        NAV TS & 100 & 30,9 & 7,2 & 4681,9 & 233,5 & 327,3 & 748,6 \\
        \hline
        LAI T & 55,7 & 31 & 10 & 1374,8 & 68,2 & 95,7 & 219,5 \\
        \hline
        LAI TS & 100 & 31,4 & 7,3 & 4873,6 & 243,1 & 340,6 & 779,2 \\
        \hline
    \end{tabular}
\end{table}

L’analyse comparative met en évidence des différences nettes entre les plots T et les plots TS. Les plots TS présentent des valeurs de $EC$, de $N$, de $P$ et de $K$ environ deux à trois fois plus élevées traduisant un enrichissement du sol en liaison avec l’amendement. Le pH égal à 10 dans les plots T est ramené à 7,2-8.2 dans les plots TS, tandis que leur humidité atteint 100\% contre 55,7-79,3\% dans les plots T. Ces résultats suggèrent ainsi une amélioration initiale de la fertilité, de la rétention en eau du sol amendé et de la régulation de l'alcalinité.

\subsection{Évolution de l'état du sol entre plots tests et plots témoins (observations)} 
Quelques jours après l’incorporation des composants de l’amendement et les premiers arrosages, des différences marquées ont été observées entre les plots témoins et les plots tests pour nos trois cultures.

Sur les plots TS, le sol est resté meuble et friable, avec une infiltration satisfaisante tout au long du cycle. Sur les plots T, en particulier NAV T et LAI T, des fissures superficielles sont apparues et le sol semblait s’affaisser sur lui-même (Figure ci-dessous). 
L’eau d’irrigation ruisselait alors vers le côté en pente du plot au lieu de s’infiltrer uniformément.
\begin{figure}[H]
    \centering
    \imgstd{images/a_fissure_chap5.jpg}
    \caption{Illustration avec une image de plot (NAV T) qui s’affaisse. }
    \label{fig:fissure5}
\end{figure}

L’eau s’infiltrait rapidement sur les plots TS, alors qu’elle nécessitait au minimum cinq minutes sur les plots T (Figures ci-dessous). 
Après des pluies nocturnes intenses, l’eau était déjà fortement infiltrée dans les plots TS, tandis que les plots T restaient dans un état plus détérioré. 
Cette différence était plus visible sur NAV T et LAI T, semés en ligne, que sur GOM T, semé en poquets. Le mode de semis et d’autres facteurs propres à chaque plot ont pu contribuer à cette différence.
\begin{figure}[H]
    \centering
    \imgstd{images/b_gombo_chap5.jpg}
    \caption{Infiltration sur les plots GOM T et GOM TS 5 minutes après irrigation observée au douzième jour. }
    \label{fig:gombo5}
\end{figure}

\begin{figure}[H]
    \centering
    \imgstd{images/c_navet_chap5.jpg}
    \caption{Infiltration sur les plots NAV T et NAV TS 5 minutes après irrigation observée au douzième jour. }
    \label{fig:navet5}
\end{figure}

Pendant la journée, le sol des plots TS restait meuble, tandis que celui des plots T \textbf{durcissait} progressivement jusqu’à devenir aussi dur que du carrelage : il supportait le poids d’une personne sans \textbf{s’effondrer et était difficile à briser à la main} (Figure suivante). 
Les plots T ont nécessité un binage plus fréquent pour faciliter l’infiltration de l’eau et l’aération du sol.
\begin{figure}[H]
    \centering
    \imgstd{images/d_sol_dur_meuble_chap5.jpg}
    \caption{Illustration avec des images de sol T dur et de sol TS meuble. }
    \label{fig:dur_meuble5}
\end{figure}

Deux autres phénomènes ont été souvent observés sur les plots T et jamais sur les plots TS (Figure ci-dessous). Après arrosage, de légères taches blanches apparaissaient à la surface du sol et devenaient plus visibles après séchage, 
sous forme d’une efflorescence blanchâtre. Par ailleurs, quelques minutes après l’irrigation, l’eau, incolore au départ, devenait noirâtre à brunâtre. Enfin, le pH mesuré sur les plots témoins avant amendement était très proche de 10.
\begin{figure}[H]
    \centering
    \imgstd{images/e_signes_chap5.jpg}
    \caption{Illustration des signes observés uniquement sur les témoins : taches blanches indiquant la présence du sel 1, toxicité 2 et tache blanche + toxicité 3. }
    \label{fig:toxique5}
\end{figure}

\subsection{Interprétation et discussion sur les observations}
Les observations précédentes sont interprétées avec prudence, car aucune analyse du sol n’a permis de les confirmer directement.

\textbf{Interprétation des symptômes} : L’ensemble des symptômes observés sur les plots T (\textbf{fissuration, croûtage de surface, dureté à l’état sec, ruissellement préférentiel, taches blanches}) est compatible avec une situation de sol sodique : 
un taux élevé de sodium échangeable provoque la dispersion des argiles, la destruction de la structure du sol, la formation d’une croûte de battance qui limite l’infiltration de l’eau, et une fissuration lors du séchage \cite{niaz_organic_2022}. 
Cette \textbf{interprétation est cohérente} avec le pH très proche de 10 mesuré avant amendement. Les taches blanches pourraient indiquer une \textbf{remontée de sel par capillarité}, sans que leur nature ait été vérifiée.

\textbf{Effet attendu de l’amendement} : L’absence d’efflorescence sur les plots TS est compatible avec un effet du biochar sur la concentration en sels en surface, par fixation du sodium et entraînement en profondeur par lessivage. 
Le sol plus meuble et perméable sur les plots TS, contre un sol durci, fissuré et sujet à l’érosion sur les plots T, est cohérent avec \textbf{l’effet structurant du compost et de la PCA}. Le binage moins fréquent sur les plots TS suggère que cet effet s’est maintenu pendant le cycle. 
L’EGB agit surtout au niveau des racines, en favorisant l’absorption du potassium au détriment du sodium. Son impact n’étant pas directement observable sur le sol, son efficacité sera jugée à travers les données de croissance et de rendement.

\textbf{Observation non expliquée} : La coloration noirâtre de l’eau après irrigation, observée uniquement sur les plots T, reste inexpliquée. 
Elle pourrait être liée à la réduction de certains éléments (fer, manganèse) ou à un excès de sodium, mais cette hypothèse n’a pas été validée analytiquement et constitue une piste pour des \textbf{travaux futurs}.

\textbf{Limites} : Les composants de l’amendement ayant été appliqués ensemble, l’effet de chacun ne peut pas être isolé. De plus, chaque culture ne compte qu’un plot par modalité, sans répétition. 
Ces résultats suggèrent donc que l’amendement a amélioré l’état du sol des plots tests dans les conditions de cet essai, sans qu’ils puissent être généralisés.

\section{Courbes d'évolution des paramètres étudiés durant l'expérimentation}
\subsection{Résultats des paramètres climatiques et hydriques et leurs relations}

\begin{figure}[H]
    \centering
    \imgstd{images/f_evolution_chap5.jpg}
    \caption{Courbes d'évolution de la température, de l’humidité, du vent et de ET0 en fonction des jours. }
    \label{fig:H_T_vent_}
\end{figure}

La température maximale moyenne se situe autour de 34,5 °C(30--39 °C, maximum observé le 10 août), tandis que la température minimale reste constante entre 25 et 27 °C, affichant une variation quotidienne moyenne de 8,6 °C. En moyenne, l'humidité relative fluctue de 57,6~\% au minimum à 91,5~\% au maximum, enregistrant jusqu'à 10 jours où elle atteint ou dépasse 95~\%. La vitesse moyenne du vent atteint 12~km/h, avec un pic enregistré à 21,3~km/h le 30 juillet. L'évapotranspiration moyenne s'élève à 5,28~mm/j, atteignant un minimum le 26/07 (avec 3,68~mm/j) et un maximum le 03/08 (avec 7,18~mm/j). Malgré des périodes humides notables, $ET_0$ varie relativement peu (5,6~mm/j pour fin juillet, 5,1 pour août et 5,2 pour septembre).

\begin{figure}[H]
    \centering
    \imgstd{images/g_kc_chap5.png}
    \caption{Evolution du Kc moyen selon la culture. }
    \label{fig:Kc_moyen5}
\end{figure}
Le navet a le Kc moyen le plus élevé (0,98) car son stade mi-saison (Kc = 1,0) est atteint plus tôt (J26) et dure plus longtemps. Le gombo et la laitue ont un Kc moyen de 0,89.

\begin{figure}[H]
    \centering
    \imgstd{images/h_ET0_ETc_chap5.png}
    \caption{Comparaison entre l'évolution de ET0 et de ETc selon la culture. }
    \label{fig:eto-etc-moyen5}
\end{figure}
 L’$ET_0$ est commune aux trois cultures (5,3 à 5,5~mm/j), puisque le climat est le même. L'ETc varie selon le $K_c$ : 4,74 (gombo), 4,61 (laitue) et 5,39~mm/j (navet). Le besoin en eau du navet est donc le plus proche de la demande climatique ($ET_c$/$ET_0$ de 0,98, contre 0,88 à 0,90 pour les deux autres).

\begin{figure}[H]
    \centering
    \imgstd{images/i_evapotranspiration_chap5.jpeg}
    \caption{Relation entre $ET_0$, $ET_c$ et la dose d’irrigation réellement apportée. }
    \label{fig:eto-etc-Di5}
\end{figure}
Avec l'apparition de la saison des pluies, $ET_0$ descend de 9,5 à 4-5~mm/j, tandis que $ET_c$ s'aligne sur le $K_c$ de chaque type de culture. Suite au 17/07, les apports d’irrigation deviennent plus rares : quasi inexistantes pour le gombo, principalement en début de croissance pour la  laitue, et de plus en plus fréquentes pour le navet.

\begin{figure}[H]
    \centering
    \imgstd{images/j_cumule_Di_chap5.jpeg}
    \caption{Relation entre la dose d’irrigation journalière et la dose d’irrigation réellement appliquée. }
    \label{fig:Di-Di-aporter5}
\end{figure}
Sur l’ensemble du cycle cultural, les doses d’irrigation réellement appliquées représentent 72~\% de la dose théorique pour le gombo, 76~\% pour la laitue et 34~\% pour le navet. Ces proportions sont calculées en rapportant la dose d’irrigation réellement appliquée à la dose théorique, puis en multipliant le résultat par 100.

\begin{figure}[H]
    \centering
    \imgstd{images/k_humidite_chap5.png}
    \caption{Evolution de l’humidité moyenne du sol selon la culture. }
    \label{fig:evolution-hum}
\end{figure}
La laitue est la plus riche en humidité, suivie par le gombo puis le navet, avec respectivement 28,2~\%, 24,6~\% et 21,6~\% d'humidité. L’ordre est cohérent avec l'irrigation constatée : le navet est le plus réclamé et le plus aride. La laitue présente également la variance la plus importante (écart-type de 10~\%, pic à 55,9~\% le 20/08). Cette élévation observée au sein du plot test de la laitue s’explique par sa position topographique au niveau de la parcelle.

\begin{figure}[H]
    \centering
    \imgstd{images/l_ph_chap5.png}
    \caption{Evolution du pH moyen du sol selon la culture. }
    \label{fig:ph-sol5}
\end{figure}
Le pH du sol est neutre à légèrement alcalin: 7,55 pour le gombo, 7,89 pour le navet, et 8,23 pour la laitue. Les divergences (approximativement 0,7 unité) demeurent en dessous de la variabilité quotidienne(écart-type approximatif de 1 unité). L'axe débute à 6, ce qui amplifie visuellement les écarts.

\begin{figure}[H]
    \centering
    \imgstd{images/m_EC_chap5.png}
    \caption{Evolution de la EC  moyenne du sol selon la culture. }
    \label{fig:ec-moyensol5}
\end{figure}
Les résultats de l'ANOVA montrent que la différence entre les moyennes de conductivité électrique des trois cultures n'est pas statistiquement significative ($F(2,\,122) = 0{,}79$ ; $p = 0{,}46$). Les moyennes observées sont de 298,4~\textmu S/cm pour la laitue ($n = 42$), 276,3~\textmu S/cm pour le gombo ($n = 55$) et 260,6~\textmu S/cm pour le navet ($n = 28$). Malgré ces différences numériques, le test ne met pas en évidence de différence significative entre les cultures au seuil de 5~\%.

\begin{figure}[H]
    \centering
    \imgstd{images/n_NPK_chap5.png}
    \caption{Evolution moyenne des teneurs en N, P, K du sol selon la culture. }
    \label{fig:evolutionNPK5}
\end{figure}
Les teneurs estimées en potassium sont les plus élevées, suivies de celles du phosphore puis de l’azote. Pour le gombo ($n = 55$), les moyennes sont de 13,35~mg/kg pour N, 18,86~mg/kg pour P et 43,74~mg/kg pour K. Pour la laitue ($n = 42$), elles sont respectivement de 14,40, 20,39 et 47,26~mg/kg. Pour le navet ($n = 28$), elles sont de 12,55, 17,73 et 41,19~mg/kg. Ces valeurs étant dérivées de la EC, elles ne constituent pas des mesures indépendantes des teneurs en $N$, $P$ et $K$. Le graphique est donc interprété de manière descriptive.

\begin{figure}[H]
    \centering
    \imgstd{images/o_boxplot_ph_chap5.png}
    \caption{Evolution de la variabilité journalière du pH du sol selon la culture. }
    \label{fig:ph-boxplot5}
\end{figure}
Le $pH$ présente une variabilité journalière dans les trois cultures, avec des valeurs comprises entre 4{,}6 et 10{,}0. Les moyennes sont de 7{,}55 pour le gombo ($n = 55$), de 8{,}23 pour la laitue ($n = 42$) et de 7{,}89 pour le navet ($n = 28$). La laitue présente ainsi le pH moyen le plus élevé, tandis que le gombo présente le plus faible. Ces différences sont décrites à partir des observations disponibles, sans conclure à leur significativité statistique en l’absence d’un test comparatif.

\begin{figure}[H]
    \centering
    \imgstd{images/p_evolution_ph_EC_chap5.png}
    \caption{Evolution temporelle du pH et de la EC du sol pour la laitue. }
    \label{fig:ph-ec-evolution5}
\end{figure}
Au cours du cycle cultural de la laitue ($n = 42$ observations journalières), la conductivité électrique moyenne augmente de 190~\textmu S/cm en juillet à 293~\textmu S/cm en août, puis à 375~\textmu S/cm en septembre, avec un maximum journalier de 898~\textmu S/cm le 8 septembre. Le $pH$ passe de 7{,}3 en juillet à environ 8{,}4 en août et en septembre, tout en présentant des fluctuations journalières et deux pics à 10. L’évolution de l’$EC$ et du $pH$ est décrite à partir des observations recueillies. Aucune conclusion statistique sur leur relation avec l’humidité n’est formulée, faute de coefficient de corrélation et de valeur $p$ disponibles.

\begin{figure}[H]
    \centering
    \imgstd{images/q_tem_hum_ph_chap5.png}
    \caption{Evolution temporelle de l'humidité, de la température et du pH du sol selon la culture. }
    \label{fig:ph-hum-temp5}
\end{figure}
Les températures moyennes du sol restent globalement comprises entre 29 et 37~°C, avec peu de différences apparentes entre les cultures. Pour la laitue, l’humidité du sol augmente d’environ 18~\% en juillet à 34~\% en septembre, une évolution observée pendant la saison des pluies. Le $pH$ présente des fluctuations au cours du cycle, avec une hausse apparente pour la laitue. Ces tendances sont descriptives et ne permettent pas, à elles seules, d’établir un lien statistique entre les paramètres étudiés.

\begin{figure}[H]
    \centering
    \imgstd{images/r_relation_ph_EC_chap5.png}
    \caption{Nuage de point mettant en relation le pH et l’EC du sol. }
    \label{fig:ph-ec-nuage5}
\end{figure}
Le nuage de points permet d’examiner visuellement la relation entre le $pH$ et la $EC$ du sol. Les coefficients de corrélation descriptifs obtenus sont de +~0,40 pour la laitue, +~0,08 pour le gombo et -~0,22 pour le navet. Ces valeurs indiquent des associations de sens et d’intensité différents selon la culture. Toutefois, en l’absence des effectifs exacts utilisés pour chaque corrélation et des résultats des tests de significativité, il est impossible de conclure à l’existence d’une relation statistiquement significative entre l'$EC$ et le $pH$.

\begin{figure}[H]
    \centering
    \imgstd{images/s_ph_EC_duSol_chap5.png}
    \caption{Relation de corrélation entre les paramètres du sol et du climat. }
    \label{fig:correlation5}
\end{figure}
Il existe une parfaite corrélation entre $EC$, l'azote, le phosphore et le potassium ($r = 1$), ce qui les rend redondants. $ET_0$ est principalement définie par l'humidité relative minimale ($r = -0{,}91$), la température maximale ($r = +0{,}80$) et le vent ($r = +0{,}63$). La corrélation entre l'humidité du sol et $EC$ est positive ($r = +0{,}54$), tandis qu'avec $ET_0$, elle est négative ($r = -0{,}49$), et l'association avec la température est faible ($r = -0{,}22$).
Les corrélations sont évaluées sur les trois cultures mélangées, et celles entre $ET_c$ et le $K_c$ (0,70) reflètent principalement la progression du stade de croissance. L'évapotranspiration est le principal moteur derrière la dynamique de l'eau. Les amendements ont un impact direct sur la salinité et la fertilité. L'efficacité du contrôle de l'irrigation nécessite l'inclusion de ces corrélations.

\subsection{Consommation d’eau : zone test contre zone témoin}

La figure ci-dessous compare le volume d’eau apporté par plot de $2\,\mathrm{m^2}$ sur le cycle de chaque culture. Le pilotage par capteur a réduit la consommation d’eau de 86~\% pour le gombo (173 L contre 1 281 L), de 76~\% pour la laitue (99 L contre 420 L) et de 66~\% pour le navet (373 L contre 1 113 L). Sur la zone témoin, l’apport est régulier, de l’ordre de 21 litres par jour, avec des plateaux lors des périodes de pluie, tandis que sur la zone test la courbe ne progresse que lorsque le capteur détecte un besoin en eau.

Cet écart est cohérent avec le principe du pilotage, qui ajuste les apports à l’état d’humidité du sol plutôt qu’à un volume fixe. La consommation plus élevée du navet parmi les plots test pourrait être liée à son développement foliaire (transpiration), sans que cela ait été vérifié. Concernant la fertilisation, l’amendement ayant été réalisé avant les cultures, la fertilisation complémentaire n’a été nécessaire que lors des pluies les plus marquées et pendant la production ; ces apports n’ont pas été chiffrés.

Ces résultats valent pour les conditions de cet essai : le témoin de la laitue a été arrêté après 20 jours (plantes mortes), l’économie de 76~\% est donc un minimum, et l’économie d’eau peut aussi refléter la meilleure rétention du sol amendé, sans qu’on puisse séparer les deux effets.

\begin{figure}[H]
    \centering
    \imgstd{images/eau_appliquee_temoin_test_comparaison.jpeg}
    \caption{Comparaison de la consommation d'eau entre test et témoin. }
    \label{fig:eauTotal}
\end{figure}

\subsection{Résultats des paramètres de croissance et de développement végétative}
\begin{figure}[H]
    \centering
    \imgstd{images/t_survivant_chap5.jpeg}
    \caption{Evolution du nombre de plantes survivantes en fonction du jour. }
    \label{fig:survivant5}
\end{figure}
La laitue témoin disparaît complètement à J20, tandis que le test lui, garde ses plants à J20, puis 5 à J40. Le gombo témoin chute de 9 à 1 plant entre J14 et J27 soit 89~\% de mortalité, pendant que le test reste stable à 4 plants de J50 à J70. Le navet témoin passe de 6 à 3 plants entre J38 et J67 or le test reste stable.

\begin{figure}[H]
    \centering
    \imgstd{images/u_nbre_feuille_chap5.png}
    \caption{Evolution moyenne du nombre de feuille en fonction du jour. }
    \label{fig:nbre-feuille5}
\end{figure}
Le gombo test atteint en moyenne  9,75 feuilles à J70, contre 5 pour le témoin, soit une hausse de 95~\%. Le navet test monte à 14,8 feuilles à J50  contre 13,0 à J67 du côté témoin (mais avec seulement 3 plants survivants). La laitue test atteint 9,8 feuilles à J40.

La chute du nombre de feuilles chez le GOM TS s’explique par l’effet de la solution carbonatée et des maladies fongiques, chez le NAV TS par les dégâts des ravageurs, chez la laitue par les fongiques.
\begin{figure}[H]
    \centering
    \imgstd{images/v_longueur_chp5.png}
    \caption{Evolution moyenne de la longueur des feuilles en fonction du jour. }
    \label{fig:longueur-feuille5}
\end{figure}
Le gombo test porte des feuilles 2 à 3 fois plus longues que le témoin au jour 35 (9,8 cm contre 4,3 cm). Le navet test atteint 21,7 cm au J50, quand le témoin n'y arrive qu'au J67, soit 31~\% d'avance. L'allongement foliaire est plus favorable dans la partie test. L'effet ressort surtout sur le gombo et le navet.
\begin{figure}[H]
    \centering
    \imgstd{images/w_largeur_chap5.png}
    \caption{Evolution moyenne de la largeur des feuilles en fonction du jour. }
    \label{fig:largeur-feuille5}
\end{figure}
Les feuilles du gombo test sont plus épaisses que celles du témoin : 11,4~cm contre 6,3~cm à J35 (+81~\%), puis 14,0~cm contre 6,0~cm à J70 (+133~\%). Le navet test atteint 8,3~cm à J50, contre 7,0~cm pour le témoin à J67. La laitue test mesure 7,2~cm à J40. 
La largeur foliaire est plus développée dans la partie test, l'effet étant plus remarqué sur le gombo en fin de période.

{\small
\begin{longtable}{|p{2.5cm}|p{1.6cm}|p{6.1cm}|p{3.6cm}|}
    \caption{Synthèse des résultats des paramètres de croissance et de développement végétatif}
    \label{tab:synthese_croissance} \\
    \hline
    \textbf{Paramètre} & \textbf{Culture} & \textbf{Résultats principaux} & \textbf{Observation} \\
    \hline
    \endfirsthead
    \hline
    \textbf{Paramètre} & \textbf{Culture} & \textbf{Résultats principaux} & \textbf{Observation} \\
    \hline
    \endhead
    Survie des plants & Laitue & Le témoin disparaît à J20, tandis que le test conserve 5 plants à J40. & Survie prolongée dans le test. \\
    \hline
    Survie des plants & Gombo & Le témoin passe de 9 à 1 plant entre J14 et J27, soit environ 89~\% de mortalité. Le test conserve 4 plants de J50 à J70. & Forte mortalité dans le témoin. \\
    \hline
    Survie des plants & Navet & Le témoin passe de 6 à 3 plants entre J38 et J67. Le test reste stable. & Effectif plus stable dans le test. \\
    \hline
    Nombre de feuilles & Gombo & À J70 : 9,75 feuilles dans le test contre 5 dans le témoin. & Écart marqué en faveur du test. \\
    \hline
    Nombre de feuilles & Navet & 14,8 feuilles dans le test à J50 contre 13 dans le témoin à J67. & Valeur supérieure atteinte plus tôt dans le test. \\
    \hline
    Nombre de feuilles & Laitue & Le test atteint 9,8 feuilles à J40. & Développement foliaire observé dans le test. \\
    \hline
    Longueur des feuilles & Gombo & À J35 : 9,8~cm dans le test contre 4,3~cm dans le témoin. & Feuilles plus longues dans le test. \\
    \hline
    Longueur des feuilles & Navet & Le test atteint 21,7~cm à J50, valeur atteinte par le témoin à J67. & Développement plus précoce dans le test. \\
    \hline
    Largeur des feuilles & Gombo & À J35 : 11,4~cm contre 6,3~cm ; à J70 : 14~cm contre 6~cm. & Écart marqué en faveur du test. \\
    \hline
    Largeur des feuilles & Navet & 8,3~cm dans le test à J50 contre 7~cm dans le témoin à J67. & Valeur supérieure atteinte plus tôt dans le test. \\
    \hline
    Largeur des feuilles & Laitue & Le test atteint 7,2~cm à J40. & Développement foliaire observé dans le test. \\
    \hline
\end{longtable}
}

Dans l'ensemble, les observations montrent des différences de survie et de développement végétatif entre les plots tests et témoins. Les écarts les plus marques concernent le gombo, notamment le nombre et la largeur des feuilles. Le navet présente aussi un développement foliaire plus précoce dans le test. Pour la laitue, la mortalité précoce des plants témoins limite la comparaison des paramètres de croissance entre les deux modalités.

Ces résultats suggèrent une évolution plus favorable des cultures dans les plots TS. Toutefois, les différences de dates d'observations et d'effectifs de plants survivants doivent être prises en compte dans l'interprétation. Les attaques parasitaires observées sur le navet en fin de cycle constituent également un facteur susceptibles d'influencer son développent.

\begin{figure}[H]
    \centering
    \imgstd{images/x_hauteur_chap5.png}
    \caption{Evolution moyenne de la hauteur des feuilles en fonction du jour. }
    \label{fig:hauteur-feuille5}
\end{figure}
Le gombo test dépasse celui du témoin  environ +96,7~\% au jour 70 (41{,}3 contre 21{,}0~cm). L'écart apparaît tôt : dès le jour 35,  21{,}3~cm(gombo test) contre 7{,}0 cm(gombo témoin), soit le triple. Même chose pour le navet, le test la longueur de feuille atteint 16{,}9~cm à J41, contre 7{,}7~cm à J38 côté témoin, environ deux fois plus. La laitue test mesure 11{,}6~cm à J40, avant la mort de deux plants sur les sept marqués. La  partie test montre donc sa supériorité très tôt, et de façon marquée, sur le gombo et le navet en particulier.

\begin{figure}[H]
    \centering
    \imgstd{images/y_box_hauteur_chap5.png}
    \caption{Distribution de la hauteur des plantes selon la culture et le traitement. }
    \label{fig:hauteur-boxplot5}
\end{figure}
Les mesures ont montré que les plantes expérimentales de gombo et de navet ont des médianes de hauteur supérieures à celles des témoins. 
Cette différence est particulièrement marquante pour le gombo, avec une médiane de 28,0 cm contre 12,0 cm pour le témoin. 
Pour le navet, la médiane est également plus haute, à 22,0 cm contre 7,5 cm. Cependant, la hauteur du vaisseau testé a été interprétée comme la longueur des feuilles, ce qui rend cette comparaison moins équivalente. La laitue témoin n'est pas présente car elle n'a pas survécu.

\begin{figure}[H]
    \centering
    \imgstd{images/z_hauteur_final_chap5.png}
    \caption{Evolution de la hauteur moyenne des plants au dernier jour. }
    \label{fig:hauteur-jourfin}
\end{figure}
On évalue d'abord la hauteur maximale sur la parcelle de gombo test (41{,}2~cm), puis sur celle de gombo témoin (21{,}0~cm),suivie de la laitue test (11{,}6~cm), du navet test (10{,}4~cm) et, finalement, du navet témoin (9{,}2~cm).
Dans la partie test, la hauteur du gombo est deux fois plus grande que celle du  témoin et une stimulation  de la croissance du navet et de la laitue.

\begin{figure}[H]
    \centering
    \imgstd{images/z1_Effectif_hauteur_chap5.png}
    \caption{Distribution des effectifs par classe de hauteur. }
    \label{fig:effectif5}
\end{figure}
La bimodalité suggère une différenciation nette entre les deux types de sol. Les hauteurs faibles se regroupent dans la partie témoin, et les hauteurs élevées dans la partie test.

\begin{figure}[H]
    \centering
    \imgstd{images/z2_comparaison_L_l_chap5.png}
    \caption{Relation entre la longueur et la largeur de la plus grande feuille. }
    \label{fig:plus-G-feuille5}
\end{figure}
La figure ci-dessus montre la relation entre la longueur et la largeur de la plus grande feuille enregistrée sur chaque plante. 
Les points se regroupent en trois nuages distincts selon la culture. Le navet se trouve dans la zone des feuilles allongées, avec une longueur supérieure à 20 cm pour une largeur comprise entre 5 et 10 cm. 
Le gombo est dans la zone des feuilles larges, où la largeur atteint environ 17 cm pour des longueurs variant de 4 à 14 cm. La laitue, quant à elle, se situe dans la zone des feuilles plus petites, surtout en largeur. 
Pour tous les points combinés, le coefficient de corrélation de Pearson est de r = 0,31, ce qui indique une faible association positive. 
Ce résultat montre que la longueur et la largeur des feuilles sont faiblement liées de manière linéaire.

\begin{figure}[H]
    \centering
    \imgstd{images/z6_gom_nav_jour11_chap5.jpg}
    \caption{Photos comparatives des NAV T \& TS et GOM T \& TS au onzième jour. }
    \label{fig:gom-nav-j11}
\end{figure}

\begin{figure}[H]
    \centering
    \imgstd{images/z9_gombo_jour22_chap5.jpg}
    \caption{Photos comparatives des GOM T \& GOM TS au vingt-deuxième jour. }
    \label{fig:gombo-j22}
\end{figure}
Cette figure ci-dessus montre la différence entre le plot GOM T et GOM TS, les cercles avec contour rouge sur le GOM T indiquent les poquets qui contiennent toujours des cultures (les gombos qui ont résisté jusqu'à ce jour malgré les conditions difficiles)

\begin{figure}[H]
    \centering
    \imgstd{images/z10_navet_jour22_chap5.jpg}
    \caption{Photos comparatives des NAV T \& NAV TS au vingt-deuxième jour. }
    \label{fig:navet-j22}
\end{figure}

\begin{figure}[H]
    \centering
    \imgstd{images/z11_gombo_jour34_chap5.jpg}
    \caption{Photos comparatives des GOM T \& GOM TS au trente-quatrième jour. }
    \label{fig:gombo-j34}
\end{figure}
La figure ci-dessus montre qu'un seul plant a survécu dans le plot témoin et dans le TS, il y a des pertes par rapport aux jours précédents. Ce qui est lié aux brûlures de la solution carbonatée.

\begin{figure}[H]
    \centering
    \imgstd{images/z12_navet_jour34_chap5.jpg}
    \caption{Photos comparatives des NAV T \& NAV TS au trente-quatrième jour. }
    \label{fig:navet-j34}
\end{figure}

\begin{figure}[H]
    \centering
    \imgstd{images/z13_laitue_jour25_chap5.jpg}
    \caption{Photos comparatives des LAI T \& LAI TS au vingt-cinquième  jour. }
    \label{fig:laitue-j25}
\end{figure}
Cette image montre non seulement la présence de plants de laitue sur le plot test,  mais aussi indique que malgré l’apport important d’eau 
lié aux fortes pluies, il y a une parfaite infiltration complète sur le test, contrairement au témoin ou l’eau reste stagnante.




















\section{Résultats phytosanitaires}
\subsection{Maladies fongiques observées et efficacité des traitements appliqués}
Les maladies fongiques observées durant l'expérimentation sont la \textbf{fusariose, le phytophthora et l’anthracnose}. Comme souligné au chapitre 3, nous avons utilisé des traitements comme la solution d’ail et la solution de bicarbonate dont les effets diffèrent.

Après utilisation de la \textbf{solution de bicarbonate}, nous avons observé des \textbf{brûlures sur les feuilles de gombo}, qui s’aggravaient sous l’effet du soleil dans la journée. Ainsi, nous avons procédé par rinçage afin de stopper ou de réduire les dégâts du bicarbonate. Par contre, \textbf{la solution d’ail} avait un \textbf{effet positif mais pas très efficace pour soigner} les plantes déjà affectées.

\begin{figure}[H]
    \centering
    \imgstd{images/z14_gombo_bruler_chap5.jpg}
    \caption{Action de la solution de bicarbonate sur les feuilles de gombo. }
    \label{fig:gombo-bruler}
\end{figure}

\section{Résultats de rendement}
\subsection{Production finale par plot}
\begin{figure}[H]
    \centering
    \imgstd{images/z15_navet_temoin_chap5.jpg}
    \caption{Rendement obtenu sur le NAV T. }
    \label{fig:rendement-navT}
\end{figure}
Sur le témoin, nous n’avons que trois plants résistants, qui ont fini par produire malgré leur retard de cycle. Cependant, les tubercules obtenus ne sont pas sains.

\begin{figure}[H]
    \centering
    \imgstd{images/z16_navet_test_chap5.jpg}
    \caption{Rendement obtenu sur le NAV TS. }
    \label{fig:rendement-navTS}
\end{figure} 
Sur le plot test, la récolte a été satisfaisante : les \textbf{tubercules étaient  sains et bien développés} (en taille et en diamètre), avec un \textbf{rendement de 4{,}8~kg} au total. C’est la première fois qu’un tel résultat est obtenu en plein champ, hors serre. \textbf{Depuis la mise en place du champ d'expérimentation Agro TIC} par nos aînés de la première ou deuxième promotion, \textbf{aucune culture de navet n’avait donné un tel résultat en plein champ}. Ce résultat soutient l’intérêt de l’approche testée.

\begin{figure}[H]
    \centering
    \imgstd{images/z17_relation_longueur_Diametre.jpeg}
    \caption{Relation entre la longueur et le diamètre de quelques tubercules de navet mesurés réalisés après récolte. }
    \label{fig:long-Diam-navet}
\end{figure} 
La croissance des tubercules sur la partie témoin est plus variable, probablement en raison de la salinité, de la compacité du sol, des attaques de parasites et du stress lié à l'eau. La production est plus homogène dans la partie test.

\begin{figure}[H]
    \centering
    \imgstd{images/z18_concurence.jpeg}
    \caption{Relation de comparaison entre quelques tubercules de navet présentant une concurrence. }
    \label{fig:navet-concurence}
\end{figure} 
La compétition entre plants est déséquilibré : la plante « Max » bénéficie d’un accès supérieur 
aux ressources au détriment de la plante « Min », les différences pouvant aller jusqu'à 50~\% en longueur tubercule et diamètre tubercule. Ceci implique une plantation trop dense.

\begin{figure}[H]
    \centering
    \imgstd{images/z19_rendement_laitue_test.jpg}
    \caption{Rendement obtenu sur la LAI TS. }
    \label{fig:rendement-laitue}
\end{figure} 
Pour la laitue, il n'y a que le test qui a donné une récolte. Les témoins n’ont pas survécu.

\begin{figure}[H]
    \centering
    \imgstd{images/z20_rendement_gombo_test.jpg}
    \caption{Rendement obtenu sur le GOM TS par intervalle d’environ 5 jours. }
    \label{fig:rendement-gombo}
\end{figure} 
Pour le gombo test la récolte est continue, pas une seule fois mais plutôt tous les 4 à 7 jours pour un nombre variant entre 7 à 10 fruits.

\subsection{Comparaison des rendements entre T et TS}

\begin{table}[H]
    \centering
    \caption{Comparaison des rendements entre plot T et plot TS}
    \label{tab:comparaison_rendement}
    \begin{tabular}{|p{2.4cm}|p{5.6cm}|p{6cm}|}
        \hline
        \textbf{Spéculation} & \textbf{Plot T} & \textbf{Plot TS} \\
        \hline
        Gombo & Aucune récolte : un plant seul a survécu sans fructification & Récolte régulière après 2 mois, presque 3 mois ; les plants survivants produisent environ 7 à 10 fruits par période de 5 à 7 jours maximum \\
        \hline
        Navet & 3 tubercules récoltés hors cycle (plus de 2 mois) avec une peau anormale & Récolte effectuée dans le cycle normal, avec un rendement de 4,8~kg sur 2~m\textsuperscript{2} \\
        \hline
        Laitue & Aucune récolte : mort totale de tous les plants & Malgré les pertes par fonte, récolte de plus d'une dizaine de pommes \\
        \hline
    \end{tabular}
\end{table}

L’écart de rendement entre plots T et TS est cohérent avec les observations précédentes : sur les plots T, la dégradation de la structure et la salinité du sol ont probablement affecté dès le départ la levée et l’installation des racines, ce qui s’accorde avec la mortalité de la laitue, le retard et la faible qualité du navet, et l’absence quasi totale de production du gombo. Sur les plots TS, le sol amendé et l’irrigation pilotée ont permis un développement jusqu’à une production effective, malgré les pertes liées aux maladies fongiques. L’amendement semble ainsi avoir conditionné non seulement l’état du sol, mais aussi la capacité des cultures à survivre et à produire.








\section{Conclusion}
Les résultats mettent en évidence des différences importantes entre les plots TS et les plots T. Sur le plan du sol, les plots TS ont conservé une structure meuble et une meilleure infiltration, tandis que les plots T ont présenté un durcissement, une fissuration, un ruissellement préférentiel et des efflorescences blanchâtres, ce qui est cohérent avec la salinité et la sodicité observées au départ. 
Les plots TS ont aussi montré une meilleure survie et une croissance plus importante (gombo TS : 41,3 cm contre 21 cm pour le témoin), qui s’est traduite par une production effective : récolte régulière pour le gombo, 4,8~kg sur $2\,\mathrm{m^2}$ pour le navet et récolte de laitue malgré les pertes, alors que les plots T ont produit très peu, voire pas du tout. Le pilotage par capteur a en outre réduit la consommation d’eau de 66~\% à 86~\% selon les cultures.

Ces observations montrent une association entre l’amendement du sol, l’amélioration de son état physique et les performances des plots TS. Certains phénomènes, notamment la coloration noirâtre de l’eau après arrosage, nécessitent toutefois des analyses complémentaires. Ces résultats constituent une base pour les recommandations et perspectives présentées dans la suite du mémoire.
\chapter*{CONCLUSION GÉNÉRALE}
\addcontentsline{toc}{chapter}{CONCLUSION GÉNÉRALE}

Ce mémoire visait à démontrer l’intérêt de l’agriculture de précision pour optimiser l’irrigation et la fertilisation en sols salins, à travers une étude comparative menée sur le champ d’expérimentation Agro TIC de l’USSEIN à Kaolack. Une zone test, conduite selon une approche intégrée (étude du sol, amendement ciblé, irrigation et fertilisation pilotées par capteur), a été comparée à une zone témoin conduite selon les pratiques traditionnelles, sur le navet, le gombo et la laitue.

Les plots tests ont conservé un sol meuble et perméable, sans efflorescences salines, avec une meilleure croissance (gombo : 41,3 cm contre 21 cm) et une production effective (navet : 4,8 kg sur 2 m², récolte de laitue), alors que les plots témoins ont donné une production très limitée, voire nulle. Ces résultats vont dans le sens de l’hypothèse H1. Le pilotage par capteur a aussi réduit la consommation d’eau de 66~\%, 76~\% et 86~\% selon les cultures, ce qui confirme l’hypothèse H2. 
Le suivi des paramètres tels que le $pH$, l'humidité et la température du sol et/ou de l'air a également permis de recueillir des informations utiles à la détection précoce d'éventuelles maladies. Ces résultats permettent de confirmer l'hypothèse H3. L'hypothèse H4 n’a pas pu être vérifiée. L’application AgriSens automatise les calculs d’évapotranspiration et de besoins en irrigation et en fertilisation.

Cette étude présente toutefois des limites : le capteur 8-en-1 estime le $NPK$ à partir de l’$EC$, qui n’indique pas la nature des ions présents ; les traitements contre les maladies se sont révélés peu efficaces, et leur détection reposait uniquement sur l’observation.

Ces limites ouvrent des perspectives telles que : utiliser un capteur donnant le $NPK$ réel, identifier précisément les ions du sol, rechercher des remèdes plus efficaces adaptés au contexte local, confirmer les maladies en laboratoire, quantifier les intrants et les coûts, tester l’adoption de l’application par les agriculteurs, et étendre l’étude à d’autres cultures.
\printbibliography[heading=bibintoc, title={BIBLIOGRAPHIE}]

\appendix
% ===== ANNEXE A : éléments retirés du chapitre 1 (section sur les calculs) =====
\chapter{DÉTAILS DES MÉTHODES DE CALCUL DES BESOINS EN EAU DES CULTURES}
\label{ann:calculs}

\section{Équation de Penman-Monteith selon la méthode FAO-56}

La méthode FAO-56 de Penman-Monteith permet de calculer l'évapotranspiration de référence ($ET_0$) à partir de données météorologiques. Elle prend en compte le rayonnement net, le flux de chaleur dans le sol, la température de l'air, la vitesse du vent et les paramètres liés à la pression de vapeur d'eau.

L'équation journalière est la suivante :

\begin{equation}
    ET_0 = \frac{0{,}408 \times \Delta \times (R_n - G) + \gamma \times \dfrac{900}{T + 273} \times u_2 \times (e_s - e_a)}{\Delta + \gamma \times (1 + 0{,}34 \times u_2)}
    \label{eq:ann_penman_monteith}
\end{equation}

Les variables de l'équation sont définies dans le Tableau~\ref{tab:ann_variables_pm}.

\begin{table}[H]
    \centering
    \caption{Variables de l'équation de Penman-Monteith}
    \label{tab:ann_variables_pm}
    \begin{tabular}{|p{2cm}|p{7.5cm}|p{3cm}|}
        \hline
        \textbf{Symbole} & \textbf{Signification} & \textbf{Unité} \\
        \hline
        $ET_0$ & Évapotranspiration de référence & mm/j \\
        \hline
        $R_n$ & Rayonnement net à la surface & MJ/m\textsuperscript{2}/j \\
        \hline
        $G$ & Flux de chaleur dans le sol & MJ/m\textsuperscript{2}/j \\
        \hline
        $T$ & Température moyenne de l'air à 2~m & °C \\
        \hline
        $u_2$ & Vitesse du vent à 2~m & m/s \\
        \hline
        $e_s$ & Pression de vapeur saturante & kPa \\
        \hline
        $e_a$ & Pression de vapeur réelle & kPa \\
        \hline
        $\Delta$ & Pente de la courbe de pression de vapeur & kPa/°C \\
        \hline
        $\gamma$ & Constante psychrométrique & kPa/°C \\
        \hline
    \end{tabular}
\end{table}

\section{Autres méthodes d'estimation de l'évapotranspiration}

Lorsque les données météorologiques nécessaires à la méthode de Penman-Monteith ne sont pas disponibles, d'autres méthodes peuvent être employées.

\begin{itemize}[label=$\bullet$]
    \item \textbf{Hargreaves-Samani :} cette méthode utilise principalement les températures minimale et maximale ainsi que le rayonnement extraterrestre. Elle est plus simple à appliquer, mais sa précision dépend des conditions climatiques et d'une éventuelle calibration locale.
    \item \textbf{Priestley-Taylor :} cette méthode repose principalement sur le bilan énergétique de la surface.
    \item \textbf{Méthode du bac d'évaporation :} elle utilise l'évaporation mesurée dans un bac, convertie en estimation de l'évapotranspiration à l'aide d'un coefficient approprié.
\end{itemize}

\section{Calcul de la réserve utile et de la réserve facilement utilisable}

La réserve utile ($RU$) correspond à la quantité d'eau que le sol peut retenir et mettre à disposition des plantes entre la capacité au champ et le point de flétrissement permanent.

Elle est calculée par :

\begin{equation}
    RU = 1000 \times (\theta_{\mathrm{CC}} - \theta_{\mathrm{PFP}}) \times Z_r
    \label{eq:ann_ru}
\end{equation}

\noindent avec :

\begin{itemize}[label=$\bullet$]
    \item $RU$ : réserve utile, en mm ;
    \item $\theta_{\mathrm{CC}}$ : teneur en eau à la capacité au champ, en m\textsuperscript{3}/m\textsuperscript{3} ;
    \item $\theta_{\mathrm{PFP}}$ : teneur en eau au point de flétrissement permanent, en m\textsuperscript{3}/m\textsuperscript{3} ;
    \item $Z_r$ : profondeur racinaire, en m.
\end{itemize}

La réserve facilement utilisable ($RFU$) représente la quantité d'eau que la culture peut prélever avant de subir un stress hydrique significatif. Elle est estimée par :

\begin{equation}
    RFU = RU \times p
    \label{eq:ann_rfu}
\end{equation}

\noindent où $p$ est la fraction de la réserve utile que la culture peut consommer avant l'apparition d'un stress hydrique significatif. Cette fraction varie selon la culture et la demande évaporative.

\section{Approche duale du coefficient cultural}

L'approche duale distingue la transpiration de la culture et l'évaporation du sol. Le coefficient cultural total est exprimé par :

\begin{equation}
    K_c = K_{cb} + K_e
    \label{eq:ann_kc}
\end{equation}

L'évapotranspiration culturale est alors calculée par :

\begin{equation}
    ET_c = (K_{cb} + K_e) \times ET_0
    \label{eq:ann_etc_dual}
\end{equation}

\noindent avec :

\begin{itemize}[label=$\bullet$]
    \item $K_{cb}$ : coefficient cultural basal, représentant principalement la transpiration de la culture ;
    \item $K_e$ : coefficient d'évaporation du sol.
\end{itemize}

Cette approche est particulièrement utile lorsque les pluies et les irrigations sont irrégulières, car l'évaporation du sol varie selon son état d'humidité.

\section{Prise en compte du stress hydrique}

Lorsque l'eau disponible dans le sol devient insuffisante, la transpiration de la culture diminue. Le coefficient de stress hydrique $K_s$, compris entre 0 et 1, permet de représenter cette réduction.

Dans l'approche duale, l'évapotranspiration culturale ajustée est estimée par :

\begin{equation}
    ET_{c,\mathrm{aj}} = (K_s \times K_{cb} + K_e) \times ET_0
    \label{eq:ann_etc_aj}
\end{equation}

Une valeur de $K_s$ proche de 1 indique une absence ou une faible limitation hydrique, tandis qu'une valeur plus faible traduit une réduction de la transpiration liée au stress hydrique.

% ===== ANNEXE B : images des détails de la confection des plots (retirées du chapitre 2) =====
\chapter{DÉTAILS DE LA CONFECTION DES PLOTS}
\label{ann:plots}

Cette annexe regroupe les photographies du labour et de la confection des plots.

\begin{figure}[H]
    \centering
    \imgstd{images/09_labour.jpg}
    \caption{Le labour. }
    \label{fig:labour2}
\end{figure}

\begin{figure}[H]
    \centering
    \imgstd{images/10_confection_planche.jpg}
    \caption{Confection de plot. }
    \label{fig:confection2}
\end{figure}

% ===== ANNEXE C : éléments retirés du chapitre 3 (partie irrigation) =====
\chapter{DÉTAIL DES CALCULS DE \texorpdfstring{$ET_0$}{ET0} ET CLASSEUR EXCEL}
\label{ann:et0}

\section{Détermination des paramètres de \texorpdfstring{$ET_0$}{ET0}}

\subsection{Détermination du rayonnement extraterrestre}

Le rayonnement extraterrestre $R_a$ ne se mesure pas, il se calcule par une formule astronomique, uniquement à partir de la latitude du lieu et du numéro du jour de l'année ($J$). Pour le déterminer, voici la méthode que nous avons adoptée \cite[chap.~3]{allen_crop_1998-1} :

\begin{equation}
    d_r = 1 + 0{,}033 \times \cos\left(2\pi \times \frac{J}{365}\right)
    \label{eq:dr}
\end{equation}
\noindent C'est la distance relative Terre-Soleil.

\begin{equation}
    \delta = 0{,}409 \times \sin\left(2\pi \times \frac{J}{365} - 1{,}39\right)
    \label{eq:declinaison}
\end{equation}
\noindent Déclinaison solaire.

\begin{equation}
    \omega_s = \arccos\left(-\tan(\text{latitude}) \times \tan(\delta)\right)
    \label{eq:angle_horaire}
\end{equation}
\noindent Angle horaire du coucher du soleil.

\begin{equation}
    R_a = \frac{24 \times 60}{\pi} \times 0{,}0820 \times d_r \times \left[\omega_s \times \sin(\text{latitude}) \times \sin(\delta) + \cos(\text{latitude}) \times \cos(\delta) \times \sin(\omega_s)\right]
    \label{eq:ra}
\end{equation}
\noindent Rayonnement extraterrestre.

Les valeurs (0,0820 ; 0,033 ; 0,409 ; 1,39 ; $24 \times 60/\pi$) sont des constantes astronomiques et solaires universelles, valables partout dans le monde.

\noindent \textbf{NB : les fonctions $\sin$, $\cos$, $\tan$ et $\arccos$ doivent être calculées avec des angles en radians.}

\subsection{Détermination du rayonnement solaire}

\begin{equation}
    R_s = K_{Rs} \times \sqrt{T_{\max} - T_{\min}} \times R_a
    \label{eq:rs_hargreaves}
\end{equation}
\noindent Formule de Hargreaves \cite{allen_crop_1998-1,koc_reference_2023}.

$K_{Rs}$ est le coefficient d'ajustement, généralement 0,16 pour les régions intérieures et 0,19 pour les régions côtières. \textbf{Dans notre cas, nous utilisons la valeur 0,16.}

\subsection{Détermination de la pression de vapeur saturante}

La pression de vapeur saturante se calcule comme suit \cite[chap.~3]{allen_crop_1998-1} :

\begin{equation}
    e_s(T_{\max}) = 0{,}6108 \times \exp\left(\frac{17{,}27 \times T_{\max}}{T_{\max} + 237{,}3}\right) \quad \text{kPa}
    \label{eq:es_tmax}
\end{equation}

\begin{equation}
    e_s(T_{\min}) = 0{,}6108 \times \exp\left(\frac{17{,}27 \times T_{\min}}{T_{\min} + 237{,}3}\right) \quad \text{kPa}
    \label{eq:es_tmin}
\end{equation}

\begin{equation}
    e_s = \frac{e_s(T_{\max}) + e_s(T_{\min})}{2}
    \label{eq:es_moy}
\end{equation}

\subsection{Détermination de la pression de vapeur réelle}

La pression de vapeur réelle se calcule comme suit \cite[chap.~3]{allen_crop_1998-1} :

\begin{equation}
    e_a = \frac{e_s(T_{\min}) \times \dfrac{RH_{\max}}{100} + e_s(T_{\max}) \times \dfrac{RH_{\min}}{100}}{2} \quad \text{kPa}
    \label{eq:ea}
\end{equation}

\subsection{Déficit de pression de vapeur}

\begin{equation}
    DPV = e_s - e_a \quad \text{kPa}
    \label{eq:dpv}
\end{equation}

\subsection{Détermination du rayonnement net}

$R_n$ dépend de $R_s$, de $T$ et de $e_a$ \cite{koc_reference_2023}.

\begin{equation}
    R_n = R_{ns} - R_{nl} \quad \text{MJ/m}^2\text{/j}
    \label{eq:rn}
\end{equation}

\noindent avec :

\begin{equation}
    R_{ns} = 0{,}77 \times R_s
    \label{eq:rns}
\end{equation}

\begin{equation}
    T = \frac{T_{\max} + T_{\min}}{2}
    \label{eq:t_moy}
\end{equation}

\noindent température moyenne journalière, utilisée plus loin dans la formule finale de $ET_0$.

\begin{equation}
    \begin{split}
        R_{nl} ={}& 4{,}903 \times 10^{-9} \times \frac{(T_{\max} + 273{,}16)^4 + (T_{\min} + 273{,}16)^4}{2} \\
        & \times \left(0{,}34 - 0{,}14\sqrt{e_a}\right) \times \left(1{,}35 \times \frac{R_s}{R_{so}} - 0{,}35\right)
    \end{split}
    \label{eq:rnl}
\end{equation}

\noindent avec $R_{so} = 0{,}75 \times R_a$, le rayonnement de ciel clair.

Il est important de préciser que, dans le calcul de $R_{nl}$, le terme de température à la puissance 4 (issu de la loi de Stefan-Boltzmann, conformément à la méthode standardisée FAO-56 \cite[chap.~3]{allen_crop_1998-1}) n'utilise pas la température moyenne $T$, mais la moyenne des puissances 4 de $T_{\max}$ et $T_{\min}$ exprimées en Kelvin. Cette nuance est distincte de la température moyenne $T = (T_{\max} + T_{\min})/2$, laquelle intervient plus loin dans le calcul de $ET_0$.

\subsection{Pente de la courbe de vapeur}

Elle se détermine de la façon suivante \cite{koc_reference_2023} :

\begin{equation}
    \Delta = \frac{4098 \times e_s}{(T + 237{,}3)^2} \quad \text{kPa/°C}
    \label{eq:pente_vapeur}
\end{equation}

\subsection{Détermination de \texorpdfstring{$G$}{G}}

$G$ représente le flux de chaleur dans le sol, c'est l'énergie utilisée pour réchauffer ou refroidir le sol. Pendant une journée, le sol reçoit de la chaleur pendant certaines heures et en restitue pendant d'autres. Donc, sur une journée, ces échanges peuvent être considérés comme faibles, voire nuls, par rapport au rayonnement net. Ainsi, si $G \approx 0$, alors $R_n - G = R_n$.

\subsection{Détermination de la constante psychrométrique}

$\gamma$ est une valeur fixe \cite{koc_reference_2023}.

\begin{equation}
    \gamma = 0{,}000665 \times P
    \label{eq:gamma}
\end{equation}

\begin{equation}
    P = 101{,}3 \times \left(\frac{293 - 0{,}0065 \times z}{293}\right)^{5{,}26}
    \label{eq:pression_atm}
\end{equation}

\noindent où $z$ est l'altitude du site, soit $z = 3$~m dans notre cas.

\begin{equation}
    P \approx 101{,}26~\text{kPa}
    \label{eq:p_valeur}
\end{equation}

\begin{equation}
    \gamma \approx 0{,}0673~\text{kPa/°C}
    \label{eq:gamma_valeur}
\end{equation}

\subsection{Vitesse du vent (\texorpdfstring{$u_2$}{u2})}

$u_2$ représente la vitesse du vent ; elle ne se calcule pas ici, elle est récupérée depuis AccuWeather.

\section{Classeur Excel de calcul automatique}

% Cette partie vient juste après la partie irrigation :
% c'est le classeur qui regroupe les formules de l'irrigation.

\begin{figure}[H]
    \centering
    \imgstd{images/19_ET0_partie2.jpg}
    \caption{Feuille de calcul automatique de $ET_0$ seconde partie. }
    \label{fig:partie2}
\end{figure}

\begin{itemize}[label=$\bullet$]
    \item \textbf{La feuille $ET_c$ :} elle récupère automatiquement la date concernée et la valeur de $ET_0$ afin de déterminer $ET_c$. Cette dernière est calculée en multipliant $ET_0$ par le $K_c$ correspondant au stade de développement de la culture. La figure ci-dessous présente un aperçu de cette feuille.
\end{itemize}

\begin{figure}[H]
    \centering
    \imgstd{images/20_ETc.jpg}
    \caption{Feuille de calcul automatique de $ET_c$. }
    \label{fig:ETC3}
\end{figure}

\begin{itemize}[label=$\bullet$]
    \item \textbf{Feuille culture :} cette feuille regroupe les $K_c$ des différentes spéculations en fonction de leurs stades de développement.
\end{itemize}

\begin{figure}[H]
    \centering
    \imgstd{images/21_culture.jpg}
    \caption{Feuille culture. }
    \label{fig:culture3}
\end{figure}

\begin{itemize}[label=$\bullet$]
    \item \textbf{Feuille Notes \& Formules :} cette feuille rassemble l'ensemble des formules utilisées pour effectuer les différents calculs.
\end{itemize}

\begin{figure}[H]
    \centering
    \imgstd{images/22_note_formule.jpg}
    \caption{Feuille Notes \& Formules. }
    \label{fig:formule3}
\end{figure}

% ===== ANNEXE D : diagrammes retirés du chapitre 4 =====
\chapter{DIAGRAMMES DE MODÉLISATION ET DE PLANIFICATION}
\label{ann:diagrammes}

\section{Diagramme de Gantt}

Cette figure ci-dessous est la représentation graphique des principales étapes chronologiques de l'ensemble de notre projet de recherche, allant de la Veille à la préparation de la soutenance.

\begin{figure}[H]
    \centering
    \includegraphics[width=\textwidth]{images/1_e_gantt_retenue.png}
    \caption{Diagramme de Gantt du projet. }
    \label{fig:gantt}
\end{figure}

% ===== ANNEXE E : éléments retirés du chapitre 5 (section 5.3) =====
\chapter{STADES PHÉNOLOGIQUES ET PHOTOS COMPARATIVES}
\label{ann:photos}

\section{Stades phénologiques clés}

\begin{figure}[H]
    \centering
    \imgstd{images/z3_stade-cles.jpg}
    \caption{Dates et durées des stades phénologiques clés pour chaque culture. }
    \label{fig:stade-cles}
\end{figure}

\section{Images comparatives des plots T et TS}

\begin{figure}[H]
    \centering
    \imgstd{images/z4_gom_nav_jour4_chap5.jpg}
    \caption{Apparition de gombo et de navet au jour 2 et jour 3 sur les plots TS (pas d'apparition sur les témoins). }
    \label{fig:gom-nav-j2}
\end{figure}

\begin{figure}[H]
    \centering
    \imgstd{images/z5_gom_nav_jour8_chap5.jpg}
    \caption{Photos comparatives des NAV T \& TS et GOM T \& TS au huitième jour. }
    \label{fig:gom-nav-j8}
\end{figure}

\begin{figure}[H]
    \centering
    \imgstd{images/z7_laitue_jour4_chap5.jpg}
    \caption{Photos comparatives de LAI T \& LAI TS au quatrième jour après repiquage. }
    \label{fig:laitue-j4}
\end{figure}

\begin{figure}[H]
    \centering
    \imgstd{images/z8_laitue_jour7_chap5.jpg}
    \caption{Photos comparatives de LAI T \& LAI TS au septième jour après repiquage. }
    \label{fig:laitue-j7}
\end{figure}

\end{document}